% ============================================================
%  Harald Høffding, Søren Kierkegaard som Filosof
%  TRANSLATION (English) — complete book
%
%  Translated from the Danish of the revised edition
%  (Copyright (c) 1919 Lindhardt og Ringhof; SAGA e-book, 2020),
%  ISBN 9788726363111. First edition: P. G. Philipsen, Copenhagen, 1892.
%
%  The Danish source text is the publisher's copyrighted modern
%  edition and is therefore NOT reproduced here; only this English
%  translation is provided. Page ranges cited in the section markers
%  are those of the 1892 first edition (per Project Runeberg) and are
%  approximate navigational aids.
%
%  Structure mirrors the book:
%    Introduction
%    I.   The Romantic-Speculative Philosophy of Religion
%    II.  Søren Kierkegaard's Older Contemporaries in Denmark
%    III. Søren Kierkegaard's Personality
%    IV.  Søren Kierkegaard's Philosophy
%           A. Epistemology            (DONE)
%           B. Ethics
%                a. The Leap
%                b. The Stages  (alpha aesthetic / beta ethical / gamma religious)
%                c. The Standard
%    V.   Søren Kierkegaard and Christianity
%           A. Personal Breakthrough
%           B. The Last Word
%    Conclusion
% ============================================================
\documentclass[12pt,a4paper]{book}
\usepackage[utf8]{inputenc}
\usepackage[T1]{fontenc}
\usepackage{libertinus}
\usepackage{libertinust1math}
\usepackage{textalpha}
\usepackage[english]{babel}
\usepackage[protrusion=true,expansion=true]{microtype}
\usepackage{setspace}
\usepackage{enumitem}
\usepackage[margin=1.2in]{geometry}
\usepackage{fancyhdr}
\usepackage{hyperref}

\hypersetup{
  pdftitle={Søren Kierkegaard as Philosopher (English translation)},
  pdfauthor={Harald Høffding},
  hidelinks
}

\pagestyle{fancy}
\fancyhf{}
\rhead{Høffding, \emph{Søren Kierkegaard as Philosopher}}
\lhead{\leftmark}
\cfoot{\thepage}
\renewcommand{\headrulewidth}{0.4pt}

\setstretch{1.3}

% Unnumbered chapter that still appears in the running header + ToC
\newcommand{\bookchapter}[1]{%
  \chapter*{#1}%
  \addcontentsline{toc}{chapter}{#1}%
  \markboth{#1}{#1}%
}

\title{\textbf{\Large Søren Kierkegaard as Philosopher}\\[10pt]
  \large Harald Høffding}
\author{}
\date{\small Translated from the Danish\\
  (revised edition, 1919; first published 1892)}

\begin{document}

\frontmatter
\maketitle
\thispagestyle{empty}

% ---- Epigraph -------------------------------------------------
\clearpage
\thispagestyle{empty}
\vspace*{0.28\textheight}
\begin{quotation}
\noindent ``What is the significance of a primitive genius? It is not so
much to bring forth something absolutely new; for there is really
nothing new under the sun; but it is to revise the universally human,
the fundamental questions. This, in the deepest sense, is honesty. And
to lack primitivity altogether---that is, the revising element---to take
everything without further ado as given custom and usage, and to let it
suffice that it is custom and usage, that is, to withdraw from
responsibility for doing likewise: that is dishonesty.''
\end{quotation}
\hfill\emph{Søren Kierkegaard's Papers}, 1847.

% ---- Translator's note ---------------------------------------
\clearpage
\thispagestyle{fancy}
\section*{Translator's note}
This is an English translation of Harald Høffding's \emph{Søren
Kierkegaard som Filosof}, first published by P.~G.~Philipsen
(Copenhagen, 1892) and here rendered from the revised edition
(\copyright~1919 Lindhardt og Ringhof; e-book, 2020). Because that
Danish edition remains under the publisher's copyright, the Danish text
is not reproduced alongside the translation; the Danish may be obtained
from the publisher (see the catalog entry for a link). Page ranges given
in the source correspond to the 1892 first edition and serve only as
navigational aids. Høffding's own notes, printed as endnotes in the
modern edition, are restored here to footnotes at their anchors. German
titles and verse are left in the original; Greek is reproduced as
printed.

\tableofcontents

\mainmatter

% ============================================================
%  INTRODUCTION  (pp. 1--4)
% ============================================================
\bookchapter{Introduction}
\label{ch:indledning}

It was a melancholy thought for Søren Kierkegaard that the time would
come when he too would fall into ``the Lecturer's, the Professor's''
hand and be made an object of exposition. Indeed, he foresaw that this
very melancholy thought would be lectured on as well. What he found
repugnant in this was surely not the prospect of being made an object of
criticism. Rather, it was something bound up very closely with his whole
cast of thought. In the first place, he detested the ``treacherous''
admiration that comes afterwards, when the struggle is over, when renown
has been won, and that whitewashes the prophet's tomb---forgetting that
it is in the relation of contemporaneity to what is great that it must
show itself what sense and will one has to acknowledge it. His life's
work was to expose the self-deception in which we so easily entangle
ourselves with respect to what he called ``the historically corked-up
Ideality.'' But in the second place it was also one of his central
thoughts that there is no continuous psychological and historical
development. To exhibit a spiritual phenomenon---a thinker's or an
author's activity---as a link in a development, as borne by definite
inner and outer conditions and as bearing and conditioning the
subsequent development: that, for him, was to drag it down, to rob it of
its value. He would not be ``set down in a Paragraph.''

If my attempt to give a characterization and an appraisal of this
greatest thinker of our people should stand in need of a defense against
his own protest, then I must first observe that I proceed from an
entirely different conception of the relation between psychology and
ethics than his. I hold the conviction that, just as those figures who
stand before us with an ideal stamp arise under quite definite
psychological and historical conditions, so their ideal
significance---insofar as it is really present---is not weakened by their
being seen in their definite limitation and conditionedness. On the
contrary, the demonstration of this is absolutely necessary if the
abiding value is to be distinguished from what is due only to transient
and purely individual circumstances. And this holds not least of a
thinker of so pronounced and singular a peculiarity as Søren Kierkegaard.
If those who are to learn from him are themselves, or wish to be,
personalities, they must first have \emph{his} personal peculiarity
determined, so that they can decide what, according to their nature and
under their circumstances, they can carry over from him. Precisely in
consequence, then, of the great weight that Søren Kierkegaard has taught
us to place on personal appropriation, on personal truth, we must
undertake a psychological-historical and a critically examining treatment
of his personality and activity.

And as for the fact that we easily come into a self-deceiving relation to
the figures of the past, I must be permitted to make the personal
confession that when I now (1890), after an interval of some twenty-odd
years, have again taken up the study of Søren Kierkegaard's writings, I
have once more had to experience the power he has to make us turn our
gaze upon ourselves, to wound us by awakening a self-knowledge we are
otherwise so inclined to avoid. There is a sting in these writings that
no one who earnestly immerses himself in them will fail to feel, whatever
presuppositions he may otherwise bring. To them applies a word from the
book of Ecclesiastes that Søren Kierkegaard often uses in his journals
and discourses: ``Keep thy foot when thou goest to the house of the
Lord!'' His interpretation of this word was: You might easily come to
learn that the ideal was higher, the demand stricter, than you had
imagined; beware, therefore, lest you come to hear the ideal!---The
presuppositions with which I now [1890] resumed the reading of these
writings were, to be sure, very different from those with which, nearly
30 years ago, as a young theological student, I began my acquaintance
with them. At that time they led me into inner struggles and conflicts,
partly of a practical-personal, partly of a theoretical kind. After long
reluctance I had to acknowledge his strict conception of the Christian as
far more justified than the usual, idyllic and harmonious one. If
Christianity is really to be one and all for us, in life as in faith,
Søren Kierkegaard has drawn the consequences of it. For a time I tried to
keep going by distinguishing between faith and knowledge and between
ideal and reality. But I learned, more through personal life-experience
than through study, that it was an impossibility to slip through in this
way. The main decision had, after all, to be whether in my personal
life, in my spiritual household, I could really say myself to be borne by
supernatural help, and whether in my way of conducting my life I could be
said to let myself be guided by the ideals and commands of religious
ethics. Self-knowledge gradually gave me, after a long time of unrest and
torment, a clear answer to this, which for me has been decisive for my
whole life. Confining myself here to what concerns the public, I will
remark that I was now led, through studies in the history of philosophy
and investigations in the domains of psychology and ethics, to trace the
conditions, forms, and laws of the human life of the spirit. I was led to
see that ``the Humane'' is by no means, as Søren Kierkegaard says in one
of his journals, merely evaporated Christianity, but that it denotes a
view of life and a conduct of life which---founded by the Greeks,
deepened and inwardized by Christianity, expanded and clarified by the
theoretical and practical labor of the modern age---cannot, to be sure,
easily be summed up in brief formulas, but which nonetheless exercises
its power even over those who combat it, just as its steady development
is not hindered by transient excesses and caricatures. It is from this
standpoint of humanism (if we must, after all, have a word ending in
``-ism'') that I now look back on Kierkegaard's activity. But for me it
has been significant to see how valuable his appearance must be said to
be even from this point of view. It is by no means with polemical intent
that I have returned to him, although several views with which I have
also otherwise collided appear in him with the clarity and consistency of
genius, which makes it especially instructive to bring them forward. But
in that, by the homage I can still bring the great thinker, I render my
thanks for what he has given me for my personal development, I will at
the same time bear witness that he can be much to many who stand far from
him in faith. In that respect he does not cease to be our
contemporary.---I hope it will be excused that I have here spoken so much
about myself. My concern was to give a single example, the one nearest to
me, of Søren Kierkegaard's influence.

% ============================================================
%  I.  THE ROMANTIC-SPECULATIVE PHILOSOPHY OF RELIGION  (pp. 5--15)
% ============================================================
\bookchapter{I.\quad The Romantic-Speculative Philosophy of Religion}
\label{ch:1}

\noindent\textbf{1.} The nineteenth century began with a union of
religion, philosophy, and art as the watchword in the world of spirit.
People cherished an enthusiastic faith that truth was one, and that
everything of value, in whatever domain and in whatever form it might
appear, was comprised in this one truth, if only one immersed oneself in
it with open sense. No barriers to thought---and yet preservation of
harmony with religious feeling and with artistic imagination: that was
the great idea that filled the leading thinkers in the first decades.
They believed that in the strict form of thought they were only
expressing what the newly awakening religious sense and the so mighty
poetic urge, under other forms, lived and breathed in. The faith of the
philosophers of religion was, in form and content, not very different
from that which Faust expresses when Gretchen is anxious about how things
stand with his religion. Faust expresses a feeling of infinity,
inwardness, and fullness, awakened by life in nature and among human
beings, and which cannot be summed up in any finite expression, nor
exhausted in any name---»Name ist Schall und Rauch!« And Gretchen thinks:

\begin{verse}
So ungefähr sagt das Pfarrer auch,\\
Nur mit ein bischen anderen Worten.
\end{verse}

\noindent These words might serve as a motto for the age's conception of
the relation between faith and knowledge, as it appears---with
characteristic and significant differences---in Schleiermacher and in
Hegel.

\medskip

\noindent\textbf{2.} Right at the century's beginning (1799)
\emph{Friedrich Schleiermacher} (b.\ 1768, d.\ 1834) published his »Reden
über die Religion an die Gebildeten unter ihren Verächtern«. He here
conceived of religion as an immediate consciousness that everything
finite has its subsistence in and through an infinite, everything
temporal in an eternal---not in such a way that the eternal and infinite
falls outside the temporal and finite, but in such a way that it is the
latter's inner essence and soul. While in our knowing and our acting we
direct ourselves toward the finite, the determinate and the limited, in
order to grasp it or to change it, in our feeling we live a whole and
undivided life, a universal life, in which opposition and limitation are
annulled, and in which the one-sidedness that clings to every theoretical
and practical striving falls away. Feeling is the domain of religion,
because it is a life in the whole and the universal---because there
announces itself in it that ``instinct for the Universe'' which is
precisely one with religion. Every feeling, according to Schleiermacher,
has a religious character when it is not morbid and corrupted. And as the
religious feeling thus expresses itself not in opposition to, but as a
deepening of, life in the world of experience, so the deity too stands
not as a being separated from this world, but as its inner unity and its
animating power. And immortality is not an existence that is to replace
the present one, but consists in this: to feel, in the midst of time, the
eternal within oneself.

Later Schleiermacher developed, in more theological form, his conception
of religion in »Der christliche Glaube« (1821). He here determines the
religious feeling as unconditioned feeling of dependence, and thereby
establishes a moment of the greatest importance for religious psychology.
It is, on the whole, Schleiermacher's fine psychological sense that gives
his works their abiding significance. Both his earlier description of the
religious feeling as universal feeling of unity, as a centralization of
the life of feeling in relation to existence, and the later determination
of it as feeling of dependence are valuable contributions to the
psychology of religion, even if more precise determinations are necessary.

Dogma is for Schleiermacher always something derived, something
secondary. Dogmatic propositions originally arose as expressions of the
religious feeling. First this feeling gives itself air in poetic and
rhetorical fashion; the forms of expression thus arising are then
developed in the form of doctrine, as confessions. And since the poetic
expressions often had to come into apparent conflict with one another,
even where the feeling that lay at their basis was in reality the same, a
reflective and dialectical treatment became necessary in order to bring
about full clarity. Through these intermediate stages the dogmatic
propositions came into being, and their validity and significance are
therefore different from such propositions as are won by a purely
theoretical or philosophical route. Schleiermacher wishes to carry
through a definite separation between the doctrine of faith and
philosophy, a separation he supports by his theory of knowledge (which we
know from the unfinished »Dialektik«, published after his death), which
critically demonstrates the limits of knowledge and holds a place open
for an immediate feeling-relation to the deity, of whose essence
theoretical thinking, in his view, can form no concluded concept.

By this fine and spirited conception Schleiermacher placed himself in
definite opposition to the two dominant tendencies of the age, orthodoxy
and rationalism. He rejected the former's external literalism and the
latter's superficial explaining-away. His work has an abiding
significance for the philosophy of religion, and he stands as the
greatest thinker theology can show beside Augustine and Zwingli.

\medskip

\noindent\textbf{3.} But Schleiermacher did not want to be merely a free
religious thinker; he wanted to present ecclesiastical theology and
believed he could do so from his standpoint. Here he was caught in a
great illusion and sacrifices at the altar of the age; for of his
personal love of truth there is, despite all that orthodox and radical
fanatics have thought, no doubt. He was the most Socratic personality in
the philosophy of that time. A rare sense for individualities and a
strong feeling of the right and significance of the individual mark his
life, as we know it through a rich and interesting collection of letters,
and animate his thinking. But he did not see that on such a standpoint an
ecclesiastical theologian cannot place himself. The Church as a
historical community rests on the assumption of revelation as objective
fact, as the absolute and adequate expression of the eternal truth, and
therefore comes forward with a claim to unconditional authority over the
feeling-life of individuals. For Schleiermacher the dogmas in reality
became only symbols, the standard subjective, and from his standpoint
there is only one step to Ludwig Feuerbach's, for whom theology is only
psychology. Symbols cannot be imposed; they must be freely chosen or
produced, and they can be changed without the underlying feeling
therefore needing to change. But the Church cannot concede this in the
case of its dogmas; for then the very concept of revelation would fall
away. For Schleiermacher even God's personality was only a symbolic
representation! But if, with respect to all religious representations, a
fundamental distinction is to be made between image and reality, then
thereby positive religion is abolished, which in any case must, on
certain points, maintain that truth has entered the world in its whole
fullness and in its adequate form.

And on closer inspection it appears at the same time that the very
determination of the religious feeling (both in its older form, as
feeling of unity, and in its later form, as feeling of dependence)
presupposes several theoretical presuppositions---more philosophy, that
is---than Schleiermacher is aware of. The religious feeling has, in his
conception, a harmonious character. It is precisely the feeling through
which all disharmony, one-sidedness, imperfection, and sin is overcome.
The feeling of dependence takes on, in the end, the character of trust
and enthusiasm. But thereby an optimistic and harmonious view of the
world is presupposed. It is presupposed that there are no experiences
that could bring about a wholly different total-mood toward existence, or
in any case give the religious feeling a less harmonious stamp. In
Schopenhauer and Søren Kierkegaard the opposition to the romantic,
harmonizing conception appears on this point as energetically as can be
conceived.

It is, on the whole, characteristic of the development in the domain of
spiritual, and especially religious, life in our century that the
oppositions grow ever sharper as time advances. The century begins with
harmony and reconciliation, but the splitting and division is consummated
with increasing emphasis. Schleiermacher did not live to see the radical
sharpening that came with Strauss and Feuerbach, whose forerunner he
himself is---partly through his critical investigations of the New
Testament writings (which led him, among other things, to reject on
historical grounds the authenticity of the first letter to Timothy and of
the accounts of Jesus' birth), partly through his psychological method.
He did, however, live to feel the sharpening on the orthodox wing. When
in the year 1821 he published for the third time his ``Discourses on
Religion to the Educated among its Despisers,'' he wrote in the preface
``that the times had changed so strikingly that the persons to whom these
discourses are addressed no longer seem to exist at all, but that when
one looked about among the educated, one would rather find it necessary
to write discourses to mopers and slaves of the letter, to superstitious
people who show both their ignorance and their unkindness by the manner
in which they condemn others.''---The times had changed, and they were to
change still more.

\medskip

\noindent\textbf{4.} Alongside Schleiermacher, and in a polemical posture
toward him, there worked at the University of Berlin in the twenties
another philosopher of religion who expresses the age's craving for
harmony between faith and knowledge to an even higher degree than he.
This was \emph{Hegel} (b.\ 1770, d.\ 1831). What distinguished this
thinker---whose writings'\footnote{Hegel's philosophy of religion is
most fully presented in the work published after his death, »Philosophie
der Religion«; but in his »Encyclopädie der philosophischen
Wissenschaften« (1817) it is already stated in more concise form. Most
easily accessible, perhaps, is his conception of Christianity in
»Philosophie der Geschichte«.} abstract form makes it difficult for most
present-day readers to penetrate to what he has of really significant in
his ideas---was his great sense for the unity of spiritual life under all
its different forms and at all its different stages. Since he now saw in
human spiritual life a form of the revelation of an eternal
world-spirit, it became possible for him to sketch an idealistic doctrine
of development, within whose framework he ranged---often with great
arbitrariness and with more or less spirited reinterpretation---everything
that his rich learning placed at his disposal from the various domains.
Since the divine expresses itself at all stages of existence and in all
forms, Hegel acknowledges no absolute oppositions. All differences and
oppositions are moments in the development, which serve to give it
fullness of content and energy, but which are themselves continually
annulled and united in a higher unity. Through the continual positing and
annulling of oppositions the development of the world proceeds. It is in
a certain sense the effect of contrast that is for Hegel the fundamental
phenomenon illuminating existence. It shows us, indeed, how the
oppositions call one another forth and cannot be separated. A higher
stage or a higher form of existence is, on the Hegelian conception,
characterized by the fact that oppositions have been run through and have
deposited their result in a new form of unity. Mediation, the
equalization and annulling of oppositions (or, as Hegel incorrectly calls
them, contradictions), became the watchword for Hegel's adherents. To
remain standing in a relation of opposition was a sign that one was
caught in the finite and external; the point was to reach the higher
unity, to unite the oppositions of existence in harmony.---An ingenious
idea and a great goal! But from the general idea to the special
theoretical and practical execution the way is long, and Hegel and his
adherents all too often confused the program with the execution.

The relation between philosophy and religion was, according to Hegel,
this: that they have the same content, only that philosophy expresses
this content in the language of abstract thought, while religion
expresses it in the language of feeling and imagination. The difference
is thus only a difference of form, which is not to touch the essential.
What the believer conceives as historical events---e.g.\ the creation of
the world, God's becoming man in Christ, and so on---are for philosophy
eternal truths and laws that hold for all times. The Christian doctrine
of the suffering god is thus for philosophy not the expression of a
single series of historical events; rather its essential content---which
only philosophy grasps in its clear, adequate form---is the truth that
the infinite and divine stirs in the finite world, while yet every finite
form is imperfect and must be shattered. It is, then, only through the
continual perishing of finite forms that the infinite unfolds itself in
the world. Only in the infinite process of ever new and ever again
perishing forms of life does the eternal life of the deity consist.---

Already from this one example it becomes clear that Hegel, just as much
as Schleiermacher, in reality makes the dogmas into symbols, only that
with him they become symbols for general ideas, not for subjective
feelings. It rested, then, with Hegel as with Schleiermacher on an
illusion when he believed he had reconciled faith and knowledge. But it
was Hegel's definite conviction that nothing of essential significance
was lost by the content of faith being transposed into philosophical
form; on the contrary, he believed that only then did it really become
the property of the human spirit. He did not see what Strauss and
Feuerbach first sharply demonstrated---that the transposition really
meant a transition to a wholly different view of life and the world than
that which is characteristic of positive religion, especially of
Christianity, of which of course there is here first and foremost
question. The whole relation between Christianity and humanity Hegel
conceived, like Schleiermacher, as a harmonious relation. Romantic
optimism is characteristic of both.

\medskip

\noindent\textbf{5.} But even purely philosophically and humanly there
was a misgiving that had to come forward in the face of the Hegelian
philosophy's fundamental thought of the unity of oppositions. Even if
there should be a higher unity in which all the oppositions of existence
find their reconciliation, we are, as existing beings, in the midst of
this world of oppositions, far from this unity. The surge of the
oppositions raises us and lowers us; we are dependent on experience, we
run up against barriers both in our thinking and in our acting, our
horizon is limited, and we often end in problematic results or run our
head against the wall; \emph{for us}, therefore, that unity is an
unattainable ideal, and the reconciliation it contains is only for the
imagination, or for faith. This essential point---the difference between
metaphysics on the one hand, psychology and ethics on the other---Hegel
overlooked, and therein lay his greatest philosophical error. The
consequence was that he, and still more his adherents, often confused the
purely abstract and logical establishing of the connection of oppositions
---e.g.\ of opposed standpoints---with an overcoming of the oppositions
through real work, or, in the ethical domain, through really having
\emph{lived} them, having practically attained a reconciliation of the
opposed tendencies between which life moves. Add to this that, so long as
life lasts, ever new oppositions, and thereby new problems, come forward;
``the higher unity'' cannot therefore be won once and for all. The theory
of the unity of oppositions thus sets up an ideal, but does not give us a
real possession. It was Hegel's great error that his endeavor, to use his
own words, aimed at ``rewriting the ideal into a system.'' Thereby he
came (as not only his philosophy of religion, but also his ethics and
theory of the state show) to violate the right both of reality and of the
ideal.

Søren Kierkegaard's opposition to Hegel acquires no small significance
through the sharp and penetrating manner in which he attacked this main
point in the Hegelian philosophy. And it must be well noted that we are
here not confronted with the ideas of a single philosopher. On the
contrary, the Hegelian doctrine on this point is really of a typical
kind. It represents---in its best meaning---a natural striving, which
human spiritual life can at no time give up, after finding connection and
harmony, after winning by its own powers mastery and survey over life.
And---in its more superficial application---it represents a
fantasy-relation to existence, an aesthetic-contemplative attitude toward
problems, crises, and oppositions, which one lets unfold as an
interesting spectacle, by whose conflicts one feels oneself untouched,
regarding oneself as the realized ``higher unity.''---Even in its best
meaning, however, this system---in which the romantic movement from the
century's beginning was clarified in abstract form---stood in need of a
serious reaction, one that in practice and theory laid an accent on
reality in its definite concrete and individual character. And in that
Søren Kierkegaard brings this reaction, as energetically as could be
wished, into the ethico-religious domain, he brings it at the same time
in a manner that can act as a standing resistance against similar
tendencies in human nature, such as later times too present.

Goethe's Faust, moreover, had already here expressed what was at stake.
When in Nostradamus' book he sees the symbol of the universe, he first
exclaims enthusiastically:

\begin{verse}
Wie Alles sich zum Ganzen webt,\\
Eins in dem Andern wirkt und lebt!
\end{verse}

\noindent---but he soon becomes melancholy that this unity of oppositions
is not given in the form of reality:

\begin{verse}
Welch Schauspiel! Aber ach! ein Schauspiel nur!
\end{verse}

\noindent Søren Kierkegaard conceived it precisely as his life's task to
impress what the preceding romantic-speculative generation had forgotten:
that life is something other and more than a spectacle for us.

% ============================================================
%  II.  SØREN KIERKEGAARD'S OLDER CONTEMPORARIES IN DENMARK  (pp. 16--27)
% ============================================================
\bookchapter{II.\quad Søren Kierkegaard's Older Contemporaries in Denmark}
\label{ch:2}

\noindent\textbf{1.} Johan Ludvig Heiberg (1791--1860) was the apostle of
Hegelian philosophy in Denmark.\footnote{A more detailed treatment of
this chapter's subject I have given later in my book \emph{Danske
Filosofer} (1909).} He first spoke up for it in a treatise from 1825,
which was a contribution to the deterministic controversy raised by
Howitz. Later he applied it especially to aesthetic and theological
subjects. He found in Hegel's philosophy an expression of the same thing
that Goethe strove for in his view of life and in his poetry. ``Hegel's
system,'' he says (\emph{Om Filosofiens Betydning for den nærværende
Tid}. Copenhagen 1833, p.~48), ``is the same as Goethe's. That two such
great minds should at the same time have arrived at the same result
cannot but awaken a favorable impression of the result\dots\ To
characterize the Hegelian philosophy in a few words, one may say that
it, like Goethean poetry, reconciles the ideal with reality, our demands
with what we possess, our wishes with what has been attained.'' As for
the relation between religion and philosophy in particular, Heiberg
expressed himself about it not only in his treatises and in his journal
\emph{Perseus}, but also in his great philosophical poems, which are
surely the finest of what our literature possesses in this direction. It
is Hegel's fundamental thought that he expresses when it is said, in the
Reformation cantata, that ``Thought ascended to the heights at the very
moment it descended into itself.'' For it is by thought's own path, by
thinking's clear elucidation of its own nature, that (according to Hegel)
the same ideas are to be found which constitute the essential content of
religion. To be sure, there is a great difference in the form: religion,
at its immediate stage, commands images and symbols, beside whose splendor
of color the forms of thought might seem poor and paltry. But Heiberg
consoles us in ``Protestantism in Nature,'' the poem in which he has
perhaps reached highest:

\begin{verse}
Turn your gaze toward your inner self!\\
What perished in your world\\
you will find within your thought.\\
All that vanished, in there,\\
gained life and existence.
\end{verse}

\noindent What has vanished is the outer, perceptible form, which the
childlike imagination misses, and with which it believes all is lost. But
for Heiberg, just as little as for Hegel, does anything go lost. He saw it
as peculiar to ``the present age'' that it wished to \emph{know} what
former ages had felt and believed; but this was, in his opinion, ``so far
from wanting to make an end of feeling and faith that, on the contrary,
it wants to consolidate them; for only when we know that they contain the
truth can we surrender ourselves to them undisturbed.''---This conception,
however, aroused misgivings among the theologians, and it was not in this
clear and resolute form that speculative philosophy of religion managed
to win wider assent here at home. It first had to undergo a dogmatic
modification of a peculiar kind.

\medskip

\noindent\textbf{2.} \emph{Hans Martensen} (1808--1884) felt himself,
already as a young man, strongly drawn to Hegel's forms of thought, which
seemed to him to open the possibility of a knowledge of the divine, while
Schleiermacher, who at the same time influenced him strongly, in the end
let the speculative problems lie. When Schleiermacher, during his visit to
Copenhagen the year before his death, was one day walking with the young
promising theological graduate, the latter asked him ``quite naively''
(see \emph{Martensens Levned} I, p.~69) whether he believed that a
knowledge of ``God's essence in itself, of the inner eternal life-process
in God'' was possible. Schleiermacher answered quite calmly: ``Ich halte
es für eine Täuschung.'' But the young theologian's speculative
enthusiasm would not be damped by the old master's warning. Hegel's forms
of thought enticed him, though he could not remain standing at Hegel, who
was too abstractly logical for him. He no doubt also sensed that, carried
through consistently, the dogmas could for Hegel, as for Schleiermacher,
become only symbols. He credited himself with the power to reach something
even greater than these thinkers had considered attainable: he wanted to
reconcile faith and knowledge in such a way that he retained the forms of
both! Characteristic of Martensen was a union of mysticism and
imaginative vision with thinking---that is, a thinking that strove to
bring about a certain coherence among the images that thrust themselves
forward before the imagination as moved by the temperament. His ideal was
the old mystics and theosophists; it was in the tracks of Meister Eckhart
and Jakob Böhme that he wished to walk. Of contemporary thinkers he most
admired the Catholic philosopher Franz Baader, who likewise wished
precisely to renew Jakob Böhme's philosophy. His distance from Hegel he
expressed early, in a critique of one of Heiberg's treatises
(\emph{Maanedsskrift for Litteratur} 1836) and in his dissertation on
``The Autonomy of Self-Consciousness in Christian Dogmatics'' (1837).
Against Hegel he objected above all to the impossibility of thinking
without presuppositions. The human being has not, after all, created
himself, but is posited by God---and thus has God as his presupposition,
his authority. In virtue of the dogma of creation, not independence
(autonomy), but dependence on God (theonomy), must become the watchword.
The task must become: speculation on the basis of faith. Earnest must be
made of the old saying that the fear of God is the beginning of wisdom.
Later (in his dogmatics) he expressed this by saying that the doctrine of
faith as a science was possible only for the consciousness reborn by faith
and grace.

One cannot deny Martensen an eye for the fact that his higher science
required a good many presuppositions; whereas it may well cause surprise
that he did not ask himself of what value a ``science'' is that rests on
so many and such enormous presuppositions. But the young theologian was
enthusiastic, and enthusiastic became those who heard him. People reckoned
``a new era'' from his appearance and expected a golden age of theology. In
his lectures on the history of recent philosophy he first led his hearers
to Hegel's philosophy as the necessary result of the development of
thought; and when they were impressed by this, there came astonishment and
admiration to the second power when he showed them the necessity and
possibility of a higher standpoint---beyond the greatest thinker of the
age! In his dogmatics he led his hearers, and later his readers, from
dogma to dogma without any decisive difficulty halting them. Trinity,
incarnation, atonement are presented as a series of images of great events
in the world of spirit, often accompanied by spirited and profound
remarks, but in such a way that every sting for thought is taken away.
That must come, presumably, from the rebirth. Only at one point does the
dogmatic thinker halt, without even his reborn consciousness being able to
clear the matter up for him: it is at the question of eternal damnation.
Here he finds---to the honor of his feeling of humanity, though not to the
honor of his orthodoxy---a stumbling-block for thought.

Philosophically, Martensen has, in the first place, offended against the
truth that thinking always needs presuppositions. He has not seen that
these presuppositions are always sought to be reduced to the least
possible, and that they are then (in the theory of knowledge) sought to be
demonstrated in their origin from human nature. We construct only those
presuppositions that are necessary in order to understand the reality
given in experience, and whose basis in human nature can be established.---
And in the second place, he has never shed any light on how it actually
comes about that the difficulties the dogmas present to the natural
consciousness fall away for the reborn consciousness. Does one then get a
different logic when one is reborn? And in that case, would it not be of
importance to have this higher logic set forth? Or should the whole
difference perhaps be that one no longer thinks? And is it, then, so easy
to cease thinking, to take reason captive? At any rate there came one who,
on this point, felt the most painful contradiction, the most passionate
collision between different forces, and who felt himself deeply outraged by
this attempt at a ``higher'' scientificity.

Seen from the religious point of view, the misgiving had to be this:
whether, after having obtained the presuppositions in question, one now
really also has the time, the desire, and the strength for speculation? For
it might be that those presuppositions were of such a kind that they had to
be acquired ever anew, that the appropriation could not take place once and
for all, and that the ever-renewed problem of appropriation made ``the high
unity'' impossible---even in the form in which Martensen, in his arrogant
attempt to outbid Hegel, had conceived it.

\medskip

\noindent\textbf{3.} It was German speculative philosophy that, with
Heiberg and Martensen, made its entry among us. But already before Søren
Kierkegaard made his mighty attack on this tendency, it had met able and
well-grounded resistance among Danish thinkers. Just as Treschow, at the
century's beginning, had by his sober critical thinking and his sound
psychology counteracted the obscurities of Schellingian philosophy, so
already in the thirties \emph{Frederik Christian Sibbern} (1785--1872) came
forward with a penetrating critique of Hegelian philosophy. Sibbern was,
by his innermost nature and by his whole development, unreceptive to a
philosophy that wished to spin everything out of itself, and that by its
speculative abstractions came into a relation of remoteness from life.
Characteristic of him in this last respect is a trait he himself relates.
``I remember that Søren Kierkegaard once, in his Hegelian period, met me on
Gammeltorv and asked me: what relation there was between philosophy and
life in reality? The question startled me, since all my philosophy aimed
at studying life and reality; but afterwards I certainly had to become
aware that for a Hegelianized thinker the question was natural, since the
Hegelians did not study philosophy existentially.''

For Sibbern, thinking always had to have a given, a factual basis.
Philosophy must be \emph{explicative}---that is, must make clear and
distinct the content presented in given experience---before it can become
\emph{speculative}, that is: set up its further-reaching hypotheses on this
basis. And to this was added that Sibbern, in the interesting philosophy of
development on which he early entered, conceived existence as a great
process, or a system of processes, proceeding from many different starting
points. It is the \emph{sporadic} character of development that he
especially emphasizes. Since the knower is himself one of these many
starting points, and is himself placed in one of these sporadic processes
of development, he cannot possibly survey the whole. Although we are links
in the universal existence, and its innermost essence must therefore stir
in us as well, still we cannot from our vantage point form a closed image
of the whole life of the world.

It was therefore firmly grounded in Sibbern's whole philosophical outlook
when he came forward with a \emph{Critique of Hegelian Philosophy} (1838),
directed especially against Heiberg---a work that contains many sound and
able thoughts and does Danish philosophy honor.---That he here called
Heiberg a philosophical dilettante caused Heiberg, some years later in
``A Soul after Death,'' to have Sibbern's writings belong to the library in
hell.

Sibbern's philosophy of religion held that the religious content should be
taken like any other given content of experience. The point was to live
oneself into it and find its value and validity. He was equally opposed to
the rationalistic or speculative transposition of living, personal faith
into abstract concept, and to the orthodox assertion of the dogmas as
outer, objective truths. It was, at first, the opposition to the aridity of
rationalism's understanding that determined Sibbern's position here. But as
the years went by, he felt himself more and more repelled by the orthodox
tendency, and he began to doubt how far the religious content could really
be called \emph{given} in the same sense as the experiences that science
takes as its basis. A definite critical attitude toward positive religious
faith comes forward in Sibbern in a series of remarkable university
programs from 1846--1849. It was not so much historical criticism that
brought him to doubt how much, in the religious domain, was really
\emph{given}, as it was a psychological conviction that what one took to be
original expressions of religious emotion could not be immediate
outpourings of personal experience, but had to be due to a more composite
development, in which manifold traditions and reflections had intervened.
He was, on the whole, led more and more away from the objective and
historical side of religion to the inner personal life, which is its
source of life. And in that he now connected with this his significant
thought that all development proceeds sporadically, from individual
starting points, he demanded that the right of subjectivity be
acknowledged. The aging thinker consistently carried through here what his
younger friend, Søren Kierkegaard, had already some years before
developed: that subjectivity is the truth---only that he does not, with
Kierkegaard, ground it merely on the necessity of personal appropriation,
but lays the weight especially on this, that it is in the individual
personalities, each for itself, that the real spiritual life stirs, and
that one would stop up these springs of life by imposing on all a common
form of view of life.

That Sibbern did not act more strongly on his contemporaries was connected
especially with his lacking the ability to shape his thoughts in an easily
accessible and condensed form. His expositions suffer from prolixity and
from various peculiarities, which stood hinderingly in the way of the
proper appreciation of the significance of this remarkable
personality.\footnote{A detailed characterization of Sibbern's personality
and of his cast of thought I have given in \emph{Mindre Arbejder}. First
Series, pp.~61--112.}

\medskip

\noindent\textbf{4.} A younger friend and colleague of Sibbern, the poet
and philosopher \emph{Poul Møller} (1794--1838), must be especially
remembered when the talk is of Søren Kierkegaard's older contemporaries,
since he was the one who stood in the closest personal relation to
Kierkegaard. Poul Møller had for a time been strongly occupied with Hegel's
philosophy, but later placed himself (as is seen from his treatise on the
possibility of proofs for immortality) critically toward it. As a
university teacher he worked most through his lectures on the history of
philosophy. He had, as is seen from his ``Stray Thoughts'' (\emph{Strøtanker}),
his eye open to the personal side of philosophizing; he kept not merely to
the objective content of the philosophical theories, but went back from
this content to the philosophizing personality. In his draft of a treatise
on affectation, Poul Møller's sense for the significance of personal truth
expresses itself especially in a characteristic way. It seems to me that
here is something that gives him an interesting position as Søren
Kierkegaard's forerunner. In reading these drafts one comes to sense what
the continual conversations between Poul Møller and Søren Kierkegaard had
turned upon.

Affectation is, according to Poul Møller, a combination of falseness and
self-deception. It arises through one's wanting to be what by one's nature
one cannot be, and therefore imagining to oneself and to others that one is
other than one really is. ``The affected person does on purpose what has no
significance except where it is called forth by the power and prompting of
nature or by a higher necessity.'' But Poul Møller has a clear eye for the
fact that it is something which, by the nature of circumstances, is not to
be avoided at certain stages of development. For when the personality is to
be developed and expanded, it must necessarily overstep its limits and take
foreign matter into itself, without at once being able to make this foreign
matter its own real property. Thereby a momentary affectation arises as a
point of transition. Affectation becomes fixed when one consistently
carries through the foreign principle, although it cannot take root in one's
personality. Moral feeling, self-regard, and immorality alike can in this
way be shells in which one hides oneself from oneself and others, taking
pleasure in them without, however, really filling them out. The use of
terminologies without insight into the corresponding thoughts, and
exaggerated symmetry in a scientific system, are more theoretical forms of
the same thing.

Poul Møller goes so far in his assertion of the importance of agreement
between the inner and the outer that he declares that ``no expression of
life has truth unless there lies in it a creative self-activity.'' He
advances this proposition with full awareness that it stands as a paradox.
For how much truth, indeed, will there be in the world if we are true only
in what we ourselves bring forth? But he holds the conviction that in the
falsely assumed, and in the objective common stamp, there is a great danger
to the originality and richness of life. ``Affectation most often comes
from one's not having the strength to fall out with the world in order to
show one's personality true\dots\ Everyone has by nature his deep stamp,
but lets it be effaced in order to please others.---He will not come
forward with his own person, does not believe in his person's infinite
depth. If each individual man, without fearing the reproach of simplicity,
judged of things just as they presented themselves to him, then glorious
phenomena would come forth.''---

On his deathbed Poul Møller sent, through Sibbern, a greeting to Søren
Kierkegaard. The three men belong together by kinship of thought. The idea
of personal truth and its significance they have all three grasped, though
only one of them was able to develop it in its whole force and rigor. And a
great sympathy did Søren Kierkegaard feel for his two teachers.---In the
introduction to the ``Posthumous Papers'' (\emph{Efterladte Papirer}) the
following is told. ``One of his somewhat younger schoolfellows has related
how he once, weighed down by heavy thoughts, met Søren Kierkegaard in
Frederiksberg Garden. `Without, moreover, any words being exchanged about
the mood that oppressed me,' the account runs, `Søren Kierkegaard suddenly
stepped up to me and asked, with a look and voice in which sympathy
trembled, whether I knew Professor Sibbern? \emph{Him} I should turn to;
\emph{that} was a \emph{whole} man; he was amiability itself, and in
intercourse with him there was something calming.' ''---As for Poul Møller,
Kierkegaard gave vent to his enthusiasm for him in the dedication to his
memory in \emph{Begrebet Angest} [The Concept of Anxiety]: ``The late
Professor Poul Martin Møller, the happy lover of Greek culture, admirer of
Homer, confidant of Socrates, interpreter of Aristotle---Denmark's joy in
the joy over Denmark, though `far traveled' always `remembered in the
Danish summer'---my admiration, my loss---to him this writing is
consecrated.''

This dedication indicates at the same time that Poul Møller introduced and
guided Kierkegaard in the study of the Greeks, which became so significant
for him. It means, on the whole, much in Kierkegaard's mouth when he calls
his late teacher the confidant of Socrates. There lies in it a hint of a
very essential spiritual connection between the two men, of which people
have not always been aware.

% ============================================================
%  III.  SØREN KIERKEGAARD'S PERSONALITY  (pp. 28--53)
% ============================================================
\bookchapter{III.\quad Søren Kierkegaard's Personality}
\label{ch:3}

\noindent\textbf{1.} There are thoughts that can grow up only in a
certain definite kind of soil and in a certain kind of air. And there is
no doubt that Søren Kierkegaard's thoughts belong among them. They are
like the wine that can grow only on volcanic ground. This ground, then,
we must first and foremost learn to know. And the ground lies deep. We
are directed back beyond himself to his family and his race, and to the
traditions that formed the spiritual atmosphere of his childhood.---For
the validity or invalidity of the thoughts, such a demonstration is not
decisive. Whatever peculiar conditions their development may be due to,
they can in themselves be equally valuable. It might well be that
especially valuable thoughts can grow up only under \emph{such}
conditions! But these conditions must be brought out and known when the
question of the thoughts' vital force under other conditions is to be
decided.---

\medskip

\noindent\textbf{2.} His family stems from West Jutland. In the West
Jutland race there is found, alongside tenacity, shrewdness, and humor,
also no small propensity to melancholy and weariness of life.

On the heath in Sæding parish near Ringkøbing, in the year 1768 a
twelve-year-old ``herd-boy'' walked with his sheep. He endured hunger and
cold and was oppressed by loneliness and forsakenness. Then in his despair
he went up onto a hill and cursed the God who had given him this wretched
life.---This boy was Søren Kierkegaard's father. In a journal entry from
1847 (\emph{Søren Kierkegaards Papirer}, edited by P.~A.~Heiberg and
V.~Kuhr, seventh volume) the following passage is found:

\begin{quote}
``The terrible thing about the man who once, as a little boy tending sheep
on the Jutland heath, suffered much evil, hungered and was famished, stood
up on a hill and cursed God---and that man was not able to forget it when
he was 82 years old.''
\end{quote}

When the first editor of Kierkegaard's papers showed Søren Kierkegaard's
brother, Bishop P.~C.~Kierkegaard, this passage, the old bishop burst into
tears and said: ``That is our father's history, and ours too!''
(H.~P.~Barfod's biography of Bishop Kierkegaard in \emph{Illustreret
Tidende} 1888). In the history of the Kierkegaard family that event plays
a great role, both as effect and as cause. It stands as an outflow of the
heavy temperament that felt everything more strongly and deeply than
others, and of the passion that drove every mood and every thought out to
its utmost consequence. But it also comes to stand as a threatening
symbol, which set the dark temperament in ever new motion and awakened the
thought of the God who visits the sins of the fathers upon the children.

The boy from the heath, Michael Pedersen Kierkegaard, was sent to
Copenhagen and apprenticed to a hosier. He worked his way up to becoming a
rich man, but never forgot that moment on the heath in which he believed
he had committed a sin against the Holy Spirit. His melancholy dread often
gave itself vent in despairing words, and it then stamped the life of the
home with its darkness, although he was otherwise a mentally lively man who
in wit and acumen could hold his own with his highly gifted sons. Already
in his fortieth year he had withdrawn from trade and lived thereafter
occupied mostly with philosophical reading; he especially studied the
German philosopher Wolff eagerly.---For the understanding of the following
characterization of Søren Kierkegaard, it must be noted that other
hereditary influences made themselves felt than those stemming from the
melancholy, brooding father. The mother is described as a nice little woman
with an even and contented temperament. A contrast that is not rare, and
that recalls the one between Goethe's parents, where likewise ``des Lebens
ernstes Führen'' derived from the father, ``die Frohnatur'' from the
mother. To be sure, the opposite elements appeared in Søren Kierkegaard in
a far more contentious and contrasting way than in that grand master of
life's harmony.

The younger son, \emph{Søren Aabye Kierkegaard}, was born in Copenhagen on
5 May 1813. Of his outward life there is not much to relate, and yet, when
one day all the posthumous papers are available, an interesting biography
of him will be able to be written---so true is it that the essential in
life, for a nature like his, is not the outward events. In Georg Brandes'
book on Søren Kierkegaard a brilliant characterization of his literary
personality is given. Something has here been accomplished that cannot be
done over. But I believe that Brandes, precisely because his eye was
directed especially toward S.~Kierkegaard's literary activity, has come to
lay \emph{too} great weight on individual events in his life, which
certainly became motives and themes in his production, but were not wholly
decisive for the direction of his spirit. And neither can I agree with
Brandes that piety and contempt were S.~Kierkegaard's two fundamental
passions. In my view they occupy a more subordinate place in his
psychology.---It must, however, be noted, especially with regard to this
last point, that when Brandes wrote his book, only the first two volumes of
the ``posthumous papers'' had appeared.\footnote{The first edition of
Kierkegaard's papers was undertaken by H.~P.~Barfod and the Swiss
theologian Gottsched. A new, critical edition is now being undertaken by
P.~A.~Heiberg and V.~Kuhr. I cite according to this edition, so far as it
is available.}

\medskip

\noindent\textbf{3.} The most decisive element for Søren Kierkegaard's
whole personality must be sought in his melancholy, as it lay in his race
and in his temperament and was nourished by his upbringing. ``I am not a
human being; I am melancholy to the very border of insanity,'' he says
(1849) in his notes, where he gives an account of his circumstances of
life. It was melancholy that lay at the basis of his great piety and of his
reverence for all authority. It brought with it a need for absolute points
of support and a shyness before the confusion of the new. He needed
something that could hold him up when everything seemed to dissolve for him
into a bottomless darkness. Only when such an absolute ground had been laid
could other sides of his nature, especially his great dialectical urge,
unfold freely and powerfully. In religious respects this led to his sharp
reflection and criticism being applied only to the question of the
appropriation of the religious content, not to the testing of the inner
coherence and validity of this content itself. And the feeling of contempt
that could stir in him is also easily understood from the consequences of
melancholy for his inner life. Accustomed as he was to walking between
abysses and terrors, ordinary, ``lightly living'' people had to make a
strange impression on him, just as melancholy itself had to lead him to
picture people's conduct and their motives in more unfavorable colors than
real experience warranted. But contempt was not an original element in his
nature. Rather there was in him a heartfelt compassion---understandable
also from the melancholy---with people, especially with those who suffer,
or whom he believed to be suffering. From several sides it is attested that
he had a remarkable eye for whether something pained those he was together
with, and that he knew how to soften and console in a manner as
considerate and delicate as no one else. He not only found comfort and
relief in going about and talking with common folk in the street, but felt
a heartfelt compassion with them and knew how to enter into their affairs.
And even at the last, when his conception of Christianity sharpened, and
when he felt the hour drawing near when the ideal was to be asserted in all
its rigor, he was for a while held back by compassion with the happiness he
was to disturb, with the unrest and suffering he was to cause where
hitherto a mild and harmonious conception of life had prevailed.

It is naturally impossible to derive the various qualities and tendencies
of a human being, especially of a richly endowed one, from a single cause.
Such people, whose significance rests precisely on their coming to make
spiritual experiences that most people are spared or miss, will have in
their nature several different starting points, several elements between
which they must struggle to bring harmony. The tasks they have to solve are
first and foremost set them by their own nature.

Alongside melancholy as the dominant element there has already been named a
power of thought, an ability for dialectical work, for dismembering
reflection. As he himself expresses it (1849): ``Bound in agonizing misery,
like a bird whose wings have been clipped, so am I, while I have yet
retained all the power of my spirit, that no doubt extraordinary power.''
Not only in reflection and thinking did this power express itself, but also
in an imagination that at every vibration, every flash of mood, could throw
out images in full splendor and with striking life, and in an art of
language such as had scarcely yet been seen in Danish literature. Favorably
situated in outward respects, he could develop these mental powers
according to his own pleasure. The use of these powers has, like the use of
every natural force, been bound up with pleasure; it was for him an
independent need to employ them, even if they were mainly employed in the
service of the prevailing moods.

In his work ``On My Work as an Author'' (\emph{Om Synspunktet for min
Forfatter\-virksomhed}) Kierkegaard attributes to himself, alongside ``the
enormous melancholy,'' an ``equally enormous facility for concealing it
under apparent cheerfulness and joy of life.'' By this remark he certainly
brings---as, on the whole, in the retrospect on his authorial activity---
more system and deliberateness into his conduct than there really was.
Alongside melancholy there lay an element of a wholly opposite nature, an
urge to give oneself over to all kinds of high-spirited play, to let
oneself go in wit and other trials of intellectual strength---an urge that
naturally connected itself with the urge to use his mental power. According
to the characterization his headmaster gave of him, after the old custom,
at his transition to the university, the childlikeness is said to have
persisted long in him; he had, it appears, a great love of freedom and
independence, and his character was lively and cheerful. The same
impression others got of him later in life. It cannot, at first, have been
dissimulation. There was in him an involuntary tendency in this direction, a
tendency that under a different upbringing might have been developed so as
to have a happy influence on his character. As it now expressed itself,
when melancholy withdrew, or as a contrast-effect of it, it could in fact
become a cover for the closed-off-ness of melancholy; but deliberately
produced it cannot have been. It was surely the same with this as with his
eager and constant roaming about the streets, which he later also presented
as part of a thought-out plan, but which at first too was an outflow of an
urge for movement, for people, and for changing impressions.

``I am a Janus bifrons---with the one face I laugh, with the other I
weep,'' he writes at 24 years of age. Neither of these faces was assumed,
even though the one came out most before others, the other in solitude. For
it is the nature of melancholy to withdraw when others are present, just as
the dwarfs vanish when the sun rises. This property of melancholy was
precisely what made it so fateful an element in Søren Kierkegaard's
personality. It was its working to make solitary, to produce isolation and
render openness toward others impossible, that became decisive for his
fate. The dull grief made him mute---but is it not a sin not to be able to
open oneself to those nearest one, not to be able to give oneself
inwardly? And do not, at any rate, the most intimate relations of communal
life become impossible where such an isolating melancholy prevails? When it
is told that Jesus drove out a devil, and it was mute---then Søren
Kierkegaard explained this as grief, melancholy, which in its egoism, in
order not to lose its dominion over the person, makes him mute, draws him
out of connection with others. And herewith there was now set for him an
enormous problem in his own inner life: was it insanity or sin that thus
possessed him? ``It is and remains the heaviest tribulation when a person
does not know whether the ground of his suffering is insanity or sin.''
Here was something that could well set the brooding dialectic in motion.
Streams of thought and series of images could develop from the one of the
possible points of view, only suddenly to vanish before new series of
representations from the other point of view---especially when there was
now added to this the emphasis produced by the strictly orthodox upbringing
and by the memory of how the father's melancholy had even led him to curse
God.

In \emph{Stages on Life's Way} the relation between father and son is
depicted: ``There was once a father and a son. A son is like a mirror in
which the father sees himself, and for the son again the father is like a
mirror in which he sees himself in the time to come. Yet they seldom
regarded each other thus, for the cheerfulness of a lively, animated
conversation was their daily intercourse. Only a few times did it happen
that the father stopped, stood before the son with a sorrowful
countenance, regarded him and said: Poor child, you are going about in a
quiet despair. Nothing further was ever said about how this was to be
understood, how true it nevertheless was.''---In the posthumous papers
Kierkegaard expresses himself far more definitely and sharply about how the
father ``unloaded all his melancholy onto a poor child,'' how his father
made him unhappy and robbed him of his childhood. ``The joy of being a
child I have indeed never had.'' And it is surely also grounded in his own
experience when he says: ``It is a rape, however well-meant it may be, to
force the child's existence into the decisive Christian categories''
(\emph{Concluding Unscientific Postscript}).

Kierkegaard's youthful years (from 22 to 25) have been illuminated by the
complete publication of his papers and by the investigations an energetic
Kierkegaard-scholar has undertaken on their basis.\footnote{\emph{P.~A.
Heiberg: En Episode i Søren Kierkegaards Ungdomsliv} (Kierkegaard Studier.
I) 1912.---\emph{Et Segment af Søren Kierkegaards religiøse Udvikling}
(Kierkegaard Studier III). 1918. (See my review of this book in
\emph{Tilskueren} 1919.)} It appears that there was a period in
Kierkegaard's youth that may be called the period of dissipation. He was
carried beyond the boundary of what he considered right and pure, and the
memory of this has ever since stood as a black spot in his life. When he
tried to explain to himself how such a fall was possible, he emphasizes the
demonically alluring power of foreboding and dread. The human being stands
before evil as before an abyss; a dizziness can seize him at the mere
thought of the possibility of casting himself out. Already in 1837 the view
of dread is expressed in the notes which he later developed in detail in
\emph{The Concept of Anxiety} (1844). It appears thus that there lies
personal experience behind the whole series of reflections on this concept
that are found in his works. And when in a note from 1837 he warns against
awakening in children, by allusions, suspicion, and untimely warnings,
representations of the possibility of evil, and thereby ``kindling the
tinder that is in every soul,'' he is surely thinking of his own upbringing
and of stray remarks of his father. When even the venerable figure he
looked up to with admiration for piety and firmness of character turned out
to suffer from inner unrest, perhaps to brood over secrets and in any case
to know terrible possibilities, this had to work on the son's mind like an
``earthquake.'' And the background of melancholy that in the father lay
behind the unrest and dread the son knew sufficiently from his own
experience; it was an inheritance that was increased by the strict
upbringing and by the father's remarks that the young man was going about
``in a quiet despair.'' It seems to have been in the year 1837 that Søren
Kierkegaard received confirmation of his foreboding as to what lay behind
the father's unrest. In his Old Testament conception the old man believed
that he had committed the sin that cannot be forgiven, because as a child
he had cursed God. Of the impression Søren Kierkegaard received of what he
had first only suspected, and then knew about the father, he has expressed
himself in a note from 1838.

\begin{quote}
``Then it was that the great earthquake occurred, the frightful upheaval
that suddenly forced upon me a new infallible law of interpretation of all
the phenomena. Then I suspected that my father's great age was not a divine
blessing, but rather a curse; that our family's excellent mental gifts
existed only to lacerate one another; then I felt the stillness of death
around me\dots\ A guilt must rest on the whole family, a punishment of God
must be over it.''
\end{quote}

Characteristic for Kierkegaard it surely is that the foreboding had for
him ``worked far more frightfully than the fact by which the truth of the
foreboding is confirmed,'' as he expresses himself in a note from January
1837. The foreboding and its confirmation put an end for him to the life in
speculation and in the mere spectator-relation in which he had an
inclination to place himself toward the world. ``The quiet despair'' could
no longer be kept down by such means.

His religious development had thereby, as by his whole temperament and his
whole upbringing, received an initial velocity that let his course run high
above the most traveled level roads. And the inner brooding to which he was
disposed now got a symbol to hold to. He had been led practically into
Old Testament religiosity in a way that falls to the lot of the fewest. The
patriarchal God-relation did not stand for him in the usual idyllic light.
And he was not among those who could make a great distinction between the
Old and the New Testament.

Søren Kierkegaard himself found his melancholy akin to the remarkable state
of feeling that has not been unusual in the monasteries of the Middle Ages,
and that was called \emph{Acedia}.\footnote{Cf.\ \emph{Den store Humor},
p.~96~f.} It expressed itself as spiritual slackening, disinclination for
religious exercises, longing for the world, and aversion to confession
(\emph{odium professionis}). He believed he found therein what his father
called the quiet despair. He interpreted the aversion to confession in a
wider sense: of a shyness about expressing oneself to others at all, and
believed he could then support the description with his own experience. He
also found that it betrayed a deep insight into human nature when the
medieval ethicists reckoned Acedia among the chief vices. It was a
monastery-sickness, and he also says occasionally: ``Had I lived in the
Middle Ages, I would probably have gone into a monastery.''---His
melancholy was in this different from the medieval Acedia, that it was far
more grounded in nature than the latter. Acedia arose, as it seems,
especially by contrast-effect, as a slackening after the strong mystical
exertion, a reaction after forced devotion---a «désenchantement de dieu,»
as a French philosopher has designated it. Yet Kierkegaard's state of
feeling surely also had this character in part. There was, in consequence
of the ``rape'' that early, by his own account, was practiced against him
(especially against those elements in his nature that pointed in a
direction other than melancholy), rich occasion for contrast-effect of the
kind mentioned.

\medskip

\noindent\textbf{4.} On the basis of melancholy---as the from the first
strongest, and by upbringing and the father's influence further potentiated
or inflamed, element of his nature---his engagement story now finds its
explanation, when we attend to the utterances in the posthumous papers.---
After he had led an aestheticizing and philosophizing youthful life, broken
several times by violent inner agitations, as at that ``earthquake'' and at
the father's death, he became engaged to a quite young girl, but broke off
the engagement the following year, to the great sorrow and offense of the
whole circle. He has depicted this episode in the third part of the
\emph{Stages}, and in the posthumous papers he returns to it again. It
seems to me that all necessary explanation lies in the often-varied
exclamation: ``Alas, she could not break the silence of my melancholy!'' He
had no happiness to share with her---on the contrary, it was first by
coming to know her bright, light temperament and immediate happiness that
he became rightly conscious of how unhappy a temperament his own was:
``Precisely her immediate youthful happiness, set beside my dreadful
melancholy and in such a relation, had to teach me to understand myself;
for how melancholy I was I had never suspected before; I really had no
standard for how happy a human being can be.''---As his high spirits and
power of thought could grow by the contrast to melancholy, so this again
grew when it got happiness as its counterpart. A web of actions and
reactions between the elements in his inner life and between his inner life
and his outer experiences, which with the power of consistency formed his
disposition in an ever more sharply pronounced style!

From what he calls his misery---the fact that ``the simplest conditions for
being a human being were denied him,'' the isolating, demonically inhibiting
melancholy---he derives not only the impossibility for him of entering into
a relation which, like marriage, rests on openness and inwardness toward
another, but also the impossibility of taking on an office. ``It has always
come easy to me to get on with people, and, accustomed as I am to being
bound, it never occurred to me that it should become so hard for me. But
here again comes my misery: I cannot, because I am not a human being,
because I am melancholy to the border of insanity---something I can indeed
gloss over so long as I am independent, but which makes me useless for a
position where I do not myself determine everything.'' Yet it was, according
to what the complete notes show, especially the black spot mentioned above
that stood for him as a hindrance, both toward the thought of marriage and
toward the thought of becoming a priest.

With admirable clarity and honesty Søren Kierkegaard gives an account of how
he came to be placed differently from other people. Had it been up to him,
he would have chosen marriage and social position, like so many others. It
was his sick temperament that made this impossible for him, while his
pecuniary circumstances on the other hand made it possible for him to live
(for a time---but it became his lifetime) without seeking paid work. It was
by no means the case that he kept from entering into any relation of
communal life on strictly ideal and principled grounds. That he became the
individual was due to the inner problems that, by his psychological
presuppositions, he got to struggle with. Only when, on
psychological-individual grounds, he was placed outside the ordinary way of
life of people, on a solitary post, did he there discover something that he
otherwise probably would not have caught sight of. What in his eyes stood as
the highest ideal showed itself to him with a clarity and a consistency
harder to reach for one who lives as a link in the manifold and conditioned
life of humanity. Thus from a lighthouse out at sea much is seen that cannot
be seen from the mainland; and yet it is not said that those who see it had
been drawn out there in order to see it. Melancholy became, so to speak, his
lookout station. Already in antiquity it was thought (it is so stated in the
so-called Aristotelian ``Problems'') that the melancholic temperament was
peculiar to geniuses. Even if by melancholy they did not then understand
quite the same as nowadays, this opinion may fit many cases. There are
experiences and discoveries made only by one who, with melancholy as a
weight beneath his feet, dives down into the sea of life.

The great personalities whose life and thought acquire decisive and typical
significance (whether now by the attraction or by the repulsion they
exert) are such as have great inner problems to solve, and who solve them
in a way that acquires value for others than themselves. When they come to
step forth actively and aggressively, the chief cause will lie in the fact
that the inner self-work needs an outer continuation. The sharp eye both for
the ideal and for what stands opposed to the ideal they have acquired in
their inner life. An inner drama has taken place before the outer one comes
to performance. What to contemporaries often looks like, or is by them
construed as, vanity, presumption, love of criticism, misanthropy---this is
in the really great the outflow of dearly bought experiences that they
cannot---without denying themselves---keep to themselves. Søren Kierkegaard
belongs decidedly to this class of personalities. That his inner life should
acquire significance for others through the clarity he could bring to some
of life's most important questions was his own conviction. In that
consisted, really, the faith he had in his mission. ``Very far back in my
recollection,'' he says (\emph{On My Work as an Author}), ``goes the thought
that in every generation there are a couple or three who are sacrificed for
the others, used to discover in frightful sufferings what comes to benefit
the others; thus I melancholically understood myself, that I was singled out
for this.''

\medskip

\noindent\textbf{5.} When Kierkegaard's melancholy, with its effects, stood
for him now as sickness, now as guilt, this surely came in part from the
fact that melancholy has a peculiar power to let the mind dwell on
possibilities it would otherwise not bring forward, and to give these
possibilities not merely a dark stamp, but a stamp of reality. This will
naturally happen all the more where there is a lively imagination and a
subtle dialectic to unfold the possibilities and work with them. Every
strong feeling has a tendency to spread in the mind and to give all emerging
representations partly an absolute character, partly a stamp of reality they
otherwise lack, just as it will also call forth recollections and
anticipations for its own explanation and clarification. (These phenomena I
have described, under the name of the expansion of feeling, in my
psychology.) Melancholy disposes to a life in possibilities---partly in
anticipation, partly in recollection. That Kierkegaard was well versed in
both is seen not only from his writings, where he both inculcates the great
importance of living in possibilities and testing oneself in them, and of
not confusing possibility and reality, but also from the posthumous papers,
which here too show how what he developed in his books was drawn from his
personal experience. With melancholy there worked here, naturally, the
poetic urge to image-formation, to the concrete development of the
representations that could serve as an expression of the moods, now a direct
expression of these, now, by contrast-effect, an expression of their
opposite. The following utterances show what power anticipation and
recollection had over him, and what danger he saw therein.

In the year 1839 (25 July) he writes in his diary: ``This is why I find so
little joy in existence: that when the thought of something awakens in my
soul, it awakens with such an energy, in so supernatural a magnitude, that I
really \emph{overstrain} myself upon it, and the ideal anticipation is for
me so far from being explanatory of existence that I rather proceed from it
too powerless to find what corresponds to the idea, too restless and, so to
speak, nerve-shaken to rest in it.''---One easily understands how melancholy
had to be increased by such constant re-living of the disproportion between
the anticipated and the real.

``For me,'' it runs in a note from about the same time, ``there is nothing
more dangerous than to recollect. Once I have recollected a life-relation,
it is impossible for me to get any interest in it again\dots\ Once I get time
to make the experience I have all too often enough made, that recollection
satiates more richly than all reality---then it is over.''---Here again it is
the life in representation that gets the precedence over the life in reality.
In the same direction points a note occasioned by a brother's death: ``I
have noticed that my grief is not one that apprehends momentarily, but one
that increases in the long run, and I am sure that when I some day grow old
I shall really come to think of the departed\dots\ As for my now deceased
brother, I am sure that the grief will first really awaken after a long
lapse of time.''

If melancholy, by its propensity to call forth and dwell in possibilities
and representations, gave the poetic urge and the dialectical drive wind in
their sails, there was thereby attained, on the other hand, a diverting
effect on melancholy; the latter was kept down by the work of thought and
imagination, which laid claim to the whole force and attention of the mind.
The whole first grand period of Kierkegaard's production (1843--46), in
which his most brilliant works fall, he himself has explained from this
point of view: ``Like that princess in the \emph{Thousand and One Nights} I
saved my life by telling stories, i.e.\ by producing. To produce was my
life. An enormous melancholy, inner sufferings of a sympathetic kind---all
this I could master when I was allowed to produce.''---Had there not been
this deep urge to work of thought and imagination, he would have gone under
spiritually.\footnote{The constantly repeated rhythmic play of melancholy
and thought-movement is depicted in the following note from 1843: ``It is
strange how strictly, in a certain sense, I am brought up. From time to time
I am set down in the dark hole; there I crawl about in torment and pain, see
nothing, no way out. Then suddenly a thought awakens in my soul as vividly
as if I had never had it before, though it is not unknown to me; but before
I have, as it were, only been wedded to it with the left hand, now I am
wedded to it with the right. When it has now settled fast in me, then I am
caressed a little, I am taken up in the arms, and I, who had shriveled up
like a grasshopper, now grow up again, healthy, thriving, glad, warm-blooded,
supple as a newborn. Thereupon I must, as it were, give my word that I will
pursue this thought to the utmost; I set my life in pawn, and now I am
harnessed in the traces. Stop I cannot, and my powers hold out. Then I am
finished, and now everything begins again.''} It was an independent element
in his nature, which under other circumstances could have been developed in
other forms than it was. It sought an outlet, and the direction of the
outlet was determined by what else stirred in his mind. But the whole
authorial activity was hardly, from the first, laid out so clearly and
definitely in one direction or with one purpose as he later believed. There
was from the first far more that was involuntary in him than he would admit
when he looked back on his activity.\footnote{The relation between the
involuntary and the deliberate in the design of Kierkegaard's production in
those three years I have discussed in \emph{Mindre Arbejder}. Second
Series, pp.~181--191.}

\medskip

\noindent\textbf{6.} If we stopped at the elements of his being hitherto
brought out, we could not expect to find in him more than---what is
naturally no small thing---a poet-nature, and perhaps also a thinker of a
psychological cast. Søren Kierkegaard has also often enough warned against
himself in this respect. When it was his conviction that people, especially
in his time, all too often confused a thought- and imagination-relation to
life and to life's great exemplars with a relation of reality, he knew, as
has been shown, this danger from himself. But he also knew the
counterweight. For just as characteristic of him as the strongly tensed life
of mood, the tireless dialectic, and the imagination obedient to every
nuance of feeling and every transition of thought, was the energetic will
with which he sought to preserve himself in concentration, to hold together
the elements of his being and lead all the streams of the rich inner life
toward one goal. He has with full right maintained that at no stage was he
merely a poet or a thinker. Already as a young man of 22 he wrote (at
Gilleleje, on a walking tour in the summer of 1835) the following:

\begin{quote}
``What I really lack is to get clear in my own mind about \emph{what I am to
do,} not about what I am to know, except insofar as a knowing must precede
every action. The point is to understand my destiny, to see what the Deity
really wants \emph{me} to do; the thing is to find a truth that is truth
\emph{for me,} to find \emph{the idea for which I am willing to live and
die.} And what use would it be to me if I discovered a so-called objective
truth\dots\ if it had \emph{no} deeper significance for \emph{myself} and for
\emph{my} life?\dots\ That was what I lacked: to lead \emph{a completely
human life} and not merely a life of knowledge, so that I should thereby
come not to base my developments of thought upon---yes, upon something one
calls objective---something that is at any rate not my own, but upon
something that is bound up with the deepest root of my existence, through
which I am, so to speak, grown into the Divine, cleave fast to it even if
the whole world should collapse. See, \emph{this is what I lack,} and
\emph{this is what I strive toward.}'' (The emphases are Kierkegaard's own.)
\end{quote}

The strong urge to think is thus from the first an urge \emph{to think in
existence,} to think in such a way that the innermost personality is
involved and is determined thereby in its core and its direction. He wants
to think in the same manner as, in one place, he writes that he has read a
book (Görres's \emph{Athanasius}): ``with my whole body---with the pit of my
heart.'' But this urge to bend everything in himself toward one purpose, to
join the separated, conflicting elements of his being together into unity,
had strong resistance to overcome. And yet the thinking of the age seemed to
have paved the way. The watchword, after all, was the higher unity, the
harmony of oppositions, and all around him he heard young speculative
virtuosos repeat these slogans in full assurance of being able to make good
on them. The thing was to mediate, to find the mediation; to remain stuck in
an either--or, in an extreme or an unreconciled relation of opposition, was a
sign of a poor head or of a narrow-minded holding-fast to a ``superseded
standpoint.'' These phrases had to grate on Søren Kierkegaard's ears when he
came from his solitary meditations or from his inner struggles, which taught
him well enough that here there was no question of a method, a knack, that
could be learned once and for all, but that it was something which, when
attempted in the reality of life---not merely in fantastic abstraction or on
paper---could demand the most strenuous willing, the most painful endurance,
and that it was a task that always had to be solved anew. But---as he wrote
in his diary in 1837---``there are many people who arrive at a life-result
the way schoolboys do; they cheat their teacher by copying the answer out of
the arithmetic book without having worked the problem themselves.''
Kierkegaard was here the strict auditor---first and foremost toward himself.
He lived through the elements of life and learned from his own experience
how hard a work it can be to construct a totality of them. A superficial
polishing-up so often covers over a botched job. He therefore declared
``the higher unity,'' with its both-and, a struggle for life and death. The
distinction, the disjunction, the either--or was given a place of honor, and
``repetition'' was set in the place of mediation. ``All talk of a higher
unity that is to unite absolute oppositions is a metaphysical assault on
ethics.'' One presupposes without more ado that everything is in order, and
that the elements join together as of themselves, so that no special act of
exertion, no new effort is necessary or must always be repeated. But often
this comes precisely from one's having no real natural ground in oneself.
For his own part he says: ``I sit and listen to the tones within me, the
glad hint of the music and the deep earnestness of the organ; to work them
together, that is a task not for a composer, but for a \emph{human being}
who, in default of greater demands on life, restricts himself to the simple
one of wanting to understand \emph{himself}.---To mediate is no art when one
has no moments in oneself.''

As life went on, the oppositions became for him stronger and stronger, the
moments more and more contentious, although his power grew at the same
time---or perhaps rather because his power grew; for had it not grown (even
to sickly overstrain), he would not have been able to hold fast such great
disharmonies and operate with them. All rational and natural coherence in
life was in the end burst asunder for him, and he passed a sentence of
condemnation on the natural life of man as corrupted and depraved in its
innermost root.

For it must indeed be well noted, and will in the following be more closely
demonstrated, that Søren Kierkegaard's indignation and polemic were most
immediately called forth by the romantic-speculative blurring of life's
oppositions---a blurring that, especially in the mouth of epigones and
parrots, became insipid and superficial---but that it was not merely such a
passing tendency of the age or of fashion that he wanted to do away with. He
aimed higher. What he wanted to do away with was really all human and
natural view of life and conduct of life, faith in life's unbroken power, in
its ability by its own natural ways to solve the problems its own unfolding
has brought with it. In itself the idea of ``the higher unity'' is also only
a modernized expression for the harmony that Greek thought in its way sought
and found, and by which no human thinking can refrain from steering, even
though it is of great importance for the richness and depth of life that one
not bargain away the oppositions, and that one not too early and on loose
ground celebrate the festivals of ``mediation.''

Kierkegaard also became ever more clearly conscious of what kind of power it
really was that he was fighting against. When the time of ``the system'' was
past, his struggle was not past. He saw clearly that the scientific view of
the world, by its naturalistic path, asserted the same continuity and inner
coherence in life that the speculative view had, in its idealistic manner,
maintained. It was the same enemy in a new shape.---Yet it was a pity that
Kierkegaard did not have before him greater representatives of humanism than
the speculative epigones. He did not have enough resistance to overcome. His
work therefore did not become \emph{as} significant as it otherwise might
have become.---

The higher unity was, according to Kierkegaard, so difficult to reach not
merely for the reason that there are sharp oppositions, conflicting moments
to overcome, but also because all human existence is in a state of
\emph{becoming}. Existence means, according to his definition expressly: to
be \emph{in time}, and on this last he lays (as we shall see more fully in
his theory of knowledge) a very great weight. So long as we are in time, we
stand ever before new possibilities and new tasks, whose solution is
problematic. All our knowledge and harmony is grounded on hindsight. ``Life
must be \emph{understood backwards}. But\dots\ it must be \emph{lived
forwards}.'' Which proposition, the more it is thought through, ends
precisely in this: that life in temporality never becomes rightly
intelligible, precisely because at no moment can I get complete rest to take
up the position backwards. Every won result must be ventured anew as a
stake, put anew up for debate, in order if possible to be won anew under the
constant development and the constant work. Here the concept of
\emph{repetition} shows itself in its significance for Kierkegaard's
thinking. It is closely connected with his emphasis on reality, on existence
as being in time. What has been won in the past has only a \emph{possible}
value in the future. The point is a conversion of this possibility into
reality; but every such conversion is a repetition, since what is to be
realized existed beforehand as a possibility, for representation. (To the
concept of repetition we shall return in connection with Søren Kierkegaard's
ethics.)---

\medskip

\noindent\textbf{7.} The relation between Søren Kierkegaard's poetic gift
and his gift as a thinker must still be touched on. The two sides in his
nature work constantly together. This places him in a position of his own in
literature, outside the usual rubrics. In a certain sense he might perhaps
have achieved a greater effect if he had been able to give poetic works that
did not require an understanding of philosophical thoughts, and works of
thought that did not present so many parentheses, elaborations, and images.
Many who would enjoy poetry by itself, or follow pure and rigorous thinking,
now pass him by. But the close cooperation of his poetizing and his thinking
sprang from the close connection of both with his own innermost nature. They
both work in the service of his personality. His urge to think was an urge
to think in connection with life-experiences and moods; therefore he
constantly needed concrete situations and examples. And the mood, which
constantly accompanies, then undertakes its little excursions, has its side
effects, to whose unfolding and clarification the ever-willing and
lightning-quick imagination lends its aid. Hence the richness and
poetry---in short turns and expressions, after parenthetical outbursts, or in
elaborating images---that is scattered about in his works.

``My thinking,'' he himself says, ``is \emph{presentic}; I have as much
imagination as dialectic.'' That he did not lose himself in abstractions
remote from life and the present he here aptly connects with the fact that
imagination and thinking work together in him. Since real experience seldom
or never presents us with a character, a passion, or a state of soul in full
and consistent development, he takes imagination to his aid and
constructs---by a thought-experiment that strictly holds to the decisive
conceptual determinations that are to be practiced---figures and situations
for the visualization and contemporization of the spiritual conditions and
problems he is occupied with. His most famous books---\emph{Either--Or},
\emph{Stages on Life's Way}, \emph{Repetition}---are such ``Attempts in
Experimenting Psychology.''---Yet I believe that in this respect too there is
more involuntariness and less systematic design in his authorial activity
than he often himself gives the appearance of. Images have surely thrust
themselves forward before their use, and he has---at any rate most
often---only afterward seen what they could be used for. And where the
constructing deliberateness has intervened more strongly, it has not
seldom been to the detriment of the poetic value, and also (as will be
sought to be shown in a later section) of the philosophical value of his
expositions.

Kierkegaard lays much weight on the need to test oneself through
possibilities, since the reality of life is not sufficient to educate the
character. But it is characteristic of him that he emphasizes only one
drawback of this method of possibilities, namely that one can easily shirk
one's way to a greater proficiency than one really possesses. ``It is
difficult to test oneself in possibility; it is as if someone, without
daring to use his voice, were to test whether he had a strong voice.
Hitherto in vain I have many times speculated on finding a means of being
able to check myself in possibility'' (\emph{Stages on Life's Way}). He does
not always bring out the other conceivable drawback: that the possibilities
constructed could be so dark and abnormal that strength and courage are used
up in the struggle with shadows, without any proficiency for taking up the
struggle of reality being fetched from this dwelling in the world of
possibilities. And moreover: with reality one can at times get finished; but
possibility is like a shadow or a fog that forms again as often as it is cut
apart; the struggle with possibilities therefore easily becomes a bailing
into the vessel of the Danaids. Kierkegaard has here worked far more as a
romantic than he wished and willed. He has for many conjured forth an
imaginary world that sucked all the strength and sap out of the real one.
There are several for whom his writings have been like the lion's den, where
all tracks led in, none back out again. Others have had to free themselves
violently from his world of possibility by one of those resolute ``jerks of
decision'' that he himself in another connection recommends.

So great was his zeal for personal truth that he expressly wanted it
established that these thought-experiments did not cover his own view.
Therefore he poetized not only characters and situations, but also authors.
The pseudonymous author-names that stand on most of his writings make up a
whole circle of individualities at different standpoints. It is his striving
after ``thinking in existence'' that leads him to these pseudonyms. What
virtuosity he possessed in thus psychologically and poetically entering into
different possible author-individualities can be seen from the little work
\emph{Prefaces} (\emph{Forord}), a collection of prefaces to various imagined
writings by imagined authors.---

His thinking is not only \emph{present,} bound to perceptible conditions and
situations; it is also, what is moreover closely connected therewith,
\emph{subjective}, that is: a thinking that proceeds from, and is consumed
in, the personal life, in the view of life and the conduct of life of the
individual, in the practical relation to life's circumstances and life's
possibilities. Such a subjective thinking is bound up with passion, for to
exist is an enormous contradiction. The subjective thinker must not swing
away from existence into abstraction; his task is rather to understand the
abstract and the universal concretely, in his existence as this individual
human being. To be the ideal in imagination or on the lecturer's platform is
not difficult; but to have to exist as an idealist is an extremely
strenuous life-task. The subjective thinker is an artist, not a man of
science, for his thinking goes together with a willing.---As such subjective
thinkers Kierkegaard especially admired Socrates and Lessing, also Jacobi
and Hamann. His doctoral dissertation dealt with Socrates (\emph{Om Begrebet
Ironi med stadigt Hensyn til Sokrates}, 1841). Both Socrates and Lessing he
has characterized in a brilliant, but also highly one-sided way; in
particular he has surely made them more ``subjective'' than they in reality
were.

It is an important idea that Kierkegaard has set up in his demand for
subjective thinking. It contains not merely the demand to leave the paths of
speculative abstraction and go back to critical reflection on the
presuppositions and barriers within which our thought has to work. But it
leads him to give a psychological and ethical introduction to a view of
life, or a theory of the art of the view of life. Søren Kierkegaard's
philosophy is such a theory. It is, to be sure, first and foremost designed
to illuminate the Christian view of life; but it is designed and carried
through in such a way that it can acquire significance for every view of
life.

Søren Kierkegaard's authorial activity falls essentially into two periods.
In the first (1843--46) he is occupied with the struggle against the
speculative-aesthetic evaporation of existence, with demonstrating
oppositions, transitional relations, and conditions in the domain of the
views of life. He here gives a kind of comparative philosophy of life. In
the second period (1847--1855) he carries on his great struggle against the
Church's mitigating conception and treatment of Christianity. This last
period too acquires significance in philosophical respects, since
Kierkegaard here draws his last consequences with regard to the relation
between religion and humanity.---In the characterization of the first period,
however, utterances from a later time may of course very well be used from
time to time, provided only that they are expressions of thoughts that were
already developed in the earlier time.

% ============================================================
%  IV.  SØREN KIERKEGAARD'S PHILOSOPHY  (pp. 54--126)
% ============================================================
\bookchapter{IV.\quad Søren Kierkegaard's Philosophy}
\label{ch:4}

The task Kierkegaard set himself as a thinker he has himself designated
as this: to make difficulties. Everything was being made easier nowadays,
he thought. In the outer world, traffic and intercourse were made easier
by railways and telegraphs; and in the world of spirit there was promised
a reconciliation of all oppositions, a harmonious unity of all conflicting
elements. So he set himself, then, to make something difficult;---for
otherwise those who wanted to make everything easy would come to lack
material in the end.

Of the men here at home by whom he was influenced during his development,
there has been talk in the foregoing. Here it shall merely be mentioned
that it was especially Martensen's lectures on speculative theology that
provoked his criticism. Already in the dissertation on Socrates there
appeared hints of a more critical, subjective, and existential conception
of thinking, and a doubt whether it should be possible to give our
knowledge as smooth and continuous a coherence as speculative philosophy
believed it had attained. When Schelling in the year 1841 began the
sensational lectures with which he wanted to set a new, higher and more
``positive'' philosophy in place of Hegel's, Kierkegaard was among the many
who hurried from all around to Berlin to become acquainted with the new
wisdom. In the beginning it seemed as if the expectations with which he
came would be fulfilled. In his diary he wrote: ``I am glad to have heard
Schelling's second hour---indescribable. I have sighed long enough, and
the thoughts have sighed within me; when he mentioned the word
`\emph{reality},' concerning philosophy's relation to reality, then the
child of thought leapt within me for joy as in Elizabeth. I remember almost
every word he said from that moment. Here perhaps clarity may come. This
one word reminded me of all my philosophical sufferings and torments.'' The
expected greater clarity failed to come; for some months later Kierkegaard
wrote to his brother: ``Dear Peter, Schelling talks the most utter
nonsense,'' and he traveled home before the lectures were at an end. Yet
that leap of the child of thought was the beginning of his subsequent life
of thought. For Schelling emphasized very strongly that speculative
philosophy could not get beyond the possible and the abstract, the
universal, and that the relation to the absolute reality---especially as
religious faith conceives it---could be posited only by an act of will,
motivated by longing and by practical, personal need: thus, as Kierkegaard
later expressed it, by a leap. The opposition between thought and
existence, between the universal and the individual, and the impossibility
of a continuous transition between them, Kierkegaard here got established
so firmly that he did not forget it. He was rightly repelled by the
arbitrary and fantastic way in which Schelling philosophized further after
having established that break in the continuity of thought, although amid
the fantastry there emerge flashes of genius. Schelling's concluding result
would hardly have appealed to him either; it aimed at healing the break
again by demonstrating a new coherence and by announcing a religiosity of
the future, ``a third kingdom'' (like Lessing earlier, and in our day
Henrik Ibsen), to which Christianity is only an introduction, and which in
its higher form is to lead beyond the opposition between mythology and
revelation through free religious self-understanding. Kierkegaard held to
the break that Schelling had established. In his philosophical writings
(especially in \emph{The Concept of Anxiety}) one can clearly notice
Schelling's influence, in part also his terminology.

A younger German philosopher to whom Kierkegaard owes much is Adolf
Trendelenburg. During his Berlin stay mentioned above he was so taken up
with Schelling and the Hegelians that he did not get to hear Trendelenburg,
which he later regretted. Trendelenburg was the first to give an
epistemological critique of the speculative method and to develop a more
critical theory of knowledge (\emph{Logische Untersuchungen}, 1840), which
has had no small influence on Kierkegaard's cast of thought. Kierkegaard
speaks of Trendelenburg as ``a man who thinks soundly and is happily formed
by the Greeks,'' and in his posthumous papers he even says: ``There is,
after all, no modern philosopher from whom I have had the benefit that I
have had from Trendelenburg.''

But whatever he owed to others, the essential ground was laid in his own
personality and in the original nature and direction of his thought.
Through solitary brooding and reflection his thoughts soon developed to
full maturity. After his return home from Berlin he lived in his studies
and his immensely productive authorial activity, from which he diverted
himself by walks, by excursions to the wooded regions of North Zealand, and
on a single occasion by a new journey to Berlin. By his roaming about the
streets and lanes he soon became a well-known Copenhagen figure, just as
Socrates was in Athens, and he too found his Aristophanes. He needed
movement for the sake of his health. But he also felt an urge to seek
relief for his melancholy and his sorrow by talking with common folk and
entering into their circumstances, helping and consoling them. He learned,
moreover, much both as a psychologist and as an author by hearing plain
people talk together: ``What one has vainly sought enlightenment about in
books, a light suddenly dawns on it by hearing one servant-girl talk with
another; an expression one has vainly tried to squeeze out of one's own
brain, sought in vain in dictionaries, even the Academy's, one hears in
passing: a common soldier says it, and he does not know what a rich man he
is. And as one who walks in the great forest, wondering at the whole,
sometimes tears off a branch, sometimes a leaf, now listens for a bird's
cry: so one goes among the crowd of people, sees now an utterance of a
state of soul, now another, learns and learns and only becomes more eager
to learn. Then one does not let oneself be deceived by books, in which the
human occurred so seldom; then one does not read about it in newspapers;
the best in the utterance, the loveliest, the little psychological trait, is
nevertheless sometimes not preserved'' (\emph{Stages on Life's Way}).---As
the deliberate everywhere comes after the involuntary, it was only later
that he conceived this roaming about the streets and lanes as a tactical
means: in order that the religious authorial activity, which was his real
task, might come to count by its own force, he wanted to weaken the
impression of his person. The technique of indirect communication required
this! We need not here enter more closely into this artifice. I have merely
wanted to sketch his circumstances of life in the time (1843--46) in which
his properly philosophical authorial activity falls.

% ------------------------------------------------------------
\section*{A.\quad Epistemology}
\label{sec:4A}
\markright{IV. Søren Kierkegaard's Philosophy --- A. Epistemology}

\noindent\textbf{1.} When Kierkegaard considers the question of what
knowledge an existing human being can attain, he is not thinking of
every possible kind of knowledge, but only of what he calls
\emph{essential} knowledge, by which he means the knowledge that bears
specifically on the knowing individual itself in its conditions of
existence---ethico-religious knowledge.\footnote{The main work here is
the \emph{Concluding Unscientific Postscript} (1846).} All other
knowledge---empirical,
mathematical, and historical---is a matter of indifference to him,
because it has nothing to do with existence. Already here he makes a
very important distinction, introduces an Either--Or. For it is surely a
great question whether much of that knowledge he regards as indifferent
is not---and has not been---of far-reaching significance for
ethico-religious knowledge. When the latter is not the same at all
times, this is due in part to that apparently ``indifferent'' knowledge,
which consciously and unconsciously comes to exert influence even on
those regions of spiritual life that one seeks to close off from it. But
Kierkegaard simply takes the correctness of that distinction as given,
and thereby marks off the boundary of his investigation.

\medskip

\noindent\textbf{2.} The train of thought in Kierkegaard's epistemology
is directed polemically both against Hegel and against those who (like
Martensen) held that once one had placed oneself on the ground of faith
and had one's consciousness reborn, one had thereby obtained the
conditions for a higher, hitherto unattainable knowledge. Against Hegel
it is observed that he has not really answered Kant's objections to a
knowledge of absolute reality. Thought cannot reach reality: for in the
very moment it would have reached it, it has converted it into
\emph{thought} reality, or into possibility; thought cannot get beyond
itself. Against speculative theology it is observed that faith has no
interest in understanding itself in any way other than by continually
remaining in the passion of faith. Of the highest things in the world
one can only be knowing as a happy or as an unhappy lover; by
approaching with speculative curiosity one in fact learns nothing. Take
away the passion that is ever renewed by uncertainty as a salutary
disciplinarian, and faith vanishes at once. The existing individual is
always in the process of becoming, and time is never concluded;
therefore full certainty can never be reached, and our knowledge can
never get beyond approximations and hypotheses, with which in the
domains of inessential knowledge we can also perfectly well content
ourselves. But if the Eternal and Unconditioned is to be grasped and
held fast, this can only happen in passion, because closure cannot be
attained. The craving for something finished and concluded is of the
Evil One, owing to indolence or impatience. We do not get beyond the
eternal striving after truth.---Kierkegaard joins with enthusiasm
Lessing's famous words: that if God held in his right hand all truth,
and in his left the ever-living striving after truth, he would grasp the
left hand, because eternal truth is not for a finite being. Lessing's
justification for this choice is, however, somewhat different from
Kierkegaard's. Lessing proceeds from the premise that it is not
possession, but only acquisition, that gives a human being worth,
inasmuch as one's powers are enlarged by seeking and one thereby grows
in perfection. This is a more positive and psychological justification
than the one we shall find in what follows from Kierkegaard, who is too
much the ascetic and too preoccupied with preparing his doctrine of the
Paradox to follow Lessing on his natural path.

More precisely, Kierkegaard's epistemology develops through the
grounding of two propositions: a logical system can be given; but a
system of existence cannot be given.

\medskip

\noindent\textbf{3.} A logical system can be given. We can develop a
coherent series of general determinations that hold for our thought. But
here one must be careful not to smuggle in concepts that are drawn from
reality. The logical system is a hypothesis whose validity for reality
is problematic. And it must be investigated in particular whether we
form this hypothesis after having come to know reality through
experience, that is, as a kind of abbreviated expression
(\emph{Abbreviatur}) for it, or whether we form it independently of
experience.---Kierkegaard expresses the wish to treat this last question
more fully on another occasion; but it never happened. It is, as one
easily sees, the question of whether the empiricist or the rationalist
conception of the fundamental presuppositions of thought is the correct
one.

A peculiar difficulty in the construction of the logical system lies,
for Kierkegaard, in how one can come to \emph{begin} it, i.e.\ get hold
of a starting point. For the simplest presuppositions are discovered
by reflection, by going back from derived assumptions to the assumptions
on which they rest---but when we have gone back far enough, how can we
halt the reflection? It has just as little reason to halt at one point
as at another, and will therefore, left to itself, continue
indefinitely. Or, as it can also be expressed: doubt cannot halt
itself. There must be a break, an arbitrary interruption. A decision
must be made to halt the regressive chain of thought and to fix the
point one has reached as the principle and starting point of the system.
But what then becomes of presuppositionlessness? For the beginning of
the logical system is then conditioned by something other than the
logical, namely by that which causes me to make the decision!

``What if, instead of talking or dreaming about an absolute beginning,
we were to speak of a \emph{leap}?''

---Against Hegel, who wanted to give an absolute system of thought that
closed upon itself like a circle, this line of reasoning by Kierkegaard
is telling. If one compares it with the superficial catechism-talk with
which Martensen ``went beyond Hegel,'' one has the difference between
them as thinkers given in a single example.---Yet it is a question
whether there is as great a difficulty about the ``beginning'' as
Kierkegaard thinks. We do indeed find our ultimate presuppositions by
reflection, or as Descartes expressed it, by doubt, or as one can
perhaps best express it, by analysis. But in our search for
presuppositions we always have a definite task in view, which is to be
solved with the help of the presuppositions we seek. The task is to
understand existence as given in experience, or some part of it. We
therefore have a definite standard for how far we need to go---namely
until we have found what is necessary for solving that task. If we can
go still further, we discover more remotely situated presuppositions
that present themselves with the same logical necessity as those first
found; but for the time being there will be no use for these
presuppositions. Perhaps advancing experience may lead us to make use of
them.---It is therefore no arbitrary act that forms the beginning of
logic, and its first principles have more than merely arbitrary validity.

\medskip

\noindent\textbf{4.} A system of existence cannot be given. This is due,
first, to the fact that all existence is in time, in the process of
becoming. The latter point especially must not be forgotten. Speculative
philosophy, which has claimed to explain everything, has forgotten that
the one who constructs the speculative system, however insignificant he
may be (``such a tiny little thingamajig as the existing Mr.\
Professor who writes the system''), still himself belongs to the
existence that is to be explained, and he is never finished with his own
existence. The system would only become possible if one could look back
on existence as concluded---but that would presuppose that one no longer
existed! As we have already heard: life is lived forwards, but
understood backwards.---And even that backward understanding, the
insight into the necessity of past development, may rest on an illusion:
for how can that which, so long as it is future, is only possible,
become necessary by belonging to the past? Should the past be more
necessary than the future? We understand the past properly only by
making ourselves contemporaneous with it---but then we see it as
becoming, not as concluded, and thus still for the most part merely as
possible! (See on this point \emph{Philosophical Fragments}.)

For the existing individual, moreover, thought and being are not one but
separate. We think either back on what has been or forward to what will
come. Simultaneity, absolute unity of thought and being, is impossible
for an existing being, and exists only as a fantasy. There can be no
continuity in our thinking, since it is continually interrupted by
absorption in existence. All unity and continuity is only abstraction.

And finally: only the individual, the concrete, exists; to exist means
to be individual. But only the universal and the abstract is the object
of thought. Herein too lies a disproportion between thought and
existence.\footnote{Before Kierkegaard, Sibbern (\emph{Hegels Filosofi
betragtet i Forhold til vor Tid}, 1838) had already criticized Hegel's
speculation from similar points of view. (See \emph{Danske Filosofer},
pp.~109--112.)}

Because a system of existence cannot be formed by our thought, it does
not follow, according to Kierkegaard, that such a system does not exist.
A system of existence is given---but only for ``one who is himself
outside existence and yet in existence---who in his eternity is
eternally concluded and yet includes existence within himself; that is,
God.'' (Since Kierkegaard defines existing as becoming, he consistently
says: God does not exist, since He is eternal.)---But for human beings,
existence is irrational, is a Paradox. The Paradox, the absurd and
contradictory, is according to Kierkegaard no concession but is, as he
puts it, an ``ontological determination,'' designating the relation
between an existing knowing spirit and the eternal truth.

It must be admitted that a Paradox has been established through the sharp
contradiction between the eternal truth (the eternal completedness of
existence in God) and the existence that is constantly in the process of
becoming. For how can, from any point of view whatsoever and for any
thought of which we can form a concept, that which is constantly
becoming stand as complete? Kierkegaard maintains firmly that ``motion
cannot be thought \emph{sub specie aeterni}'' [i.e.\ from the point of
view of eternity, apart from time]; but how then can God be knowing
about existence, about the whole of existence in its process of
becoming, since indeed ``existence cannot be thought without motion''?
---Here two thoughts have been violently yoked together, both of which
Kierkegaard demands be held fast. Epistemology must surely here turn his
favourite determination, Either--Or, against him and say: either the
concept of the absolute system in God must be an illusion, or else the
continual becoming must be a mere appearance.---Kierkegaard would not
need, if he chose the first alternative, to give up his theistic
presuppositions. Like other philosophical theists he could draw the
consequence that God himself must be in the process of becoming, that
God's being is not concluded. This idea was developed simultaneously
with Kierkegaard's philosophical production by the true founder of
modern philosophical theism, C.~H.~Weisse. Thereby the sharp antithesis
would be resolved---but the absolutely concluded element in the idea of
God would admittedly be sacrificed at the same time.---If one holds to
the view that the existence we know unfolds in time, and that we cannot
form any conception whatsoever of any other form of existence, then
there is no question of a concluded system. But existence then becomes
no Paradox either. It becomes irrational in the mathematical sense,
inasmuch as only an approximate knowledge of it becomes possible. It
stands before us like $\sqrt{2}$, which we can continually determine
with more decimal places without exhausting it. Existence could be so
rich and many-sided that there would be material enough in it for the
most energetic and comprehensive thought we can conceive. But that is
something entirely different from its being paradoxical, that is to say,
from its presenting us with insurmountable contradictions. The
inexhaustibility of existence for thought corresponds exactly to
Lessing's idea of the eternal striving after truth. If we hold fast to
this idea, we must as its objective counterpart assume a continually
becoming, and in its becoming unsurveillable, existence.\footnote{What is
here hinted at, I have developed in detail in \emph{Den menneskelige
Tanke}, pp.~306--375.}

Kierkegaard proceeds from the premise that we cannot, even with respect
to ``essential'' knowledge, content ourselves with approximations and
hypotheses. But why ever not? Provided they are genuine approximations,
and provided they are hypotheses that are increasingly verified by
experience? Then we have reason enough not to abandon our hope of
gaining clearer understanding and of being able practically to hold our
own in the struggle with the world. And it is such a hope that sustains
all our drive for knowledge.---But it will become apparent in what
follows that Kierkegaard treats the concept of approximation in a
remarkable way.

\medskip

\noindent\textbf{5.} That the ``essential'' truth is paradoxical is
grounded by Kierkegaard in the fact that the thought of God as the
absolutely eternal must be held fast in the midst of the unceasing
becoming of existence. On this depends also the fact that essential truth
can be held fast only with the passion of faith.

If one asks how the idea of God arises for Kierkegaard, his answer is
that God is a postulate, a last resort, without which we cannot endure
to live. An image that would aptly express his thought would be the
following: God is an anchor of distress that we need on the wild sea
of existence, and at which, amid the surging waves, there is a constant
tugging. We need an absolutely fixed point of support, a hook on which
the chain of existence can be fastened, a point where the infinite
reflection can find rest. But what is postulated is qualitatively
different from the postulating being: not even the concept of existence
is supposed to apply to God! God does not, according to Kierkegaard,
stand in relation to the human being as a superlative ideal. ``Between
God and man there is an absolute difference''; ``the qualities are
absolutely different.''---But then the situation is such that the last
resort to which the human being in existence anxiously flies itself
creates new difficulties and increases the suffering: for now the task
is to unite the absolutely different qualities, to think them together!
---Faith in God arises when the human being ``with infinite interest
inquires after a reality other than its own''---but thereby new problems
are immediately posed.

One understands that Kierkegaard, given his conception, could not have
it as easy as Martensen in building speculative theology on the ground
of faith. This ground was of a volcanic nature, was in constant motion,
just as it itself owed its origin to a revolution. Here there was
neither time nor place to build, but only to struggle anxiously to
preserve what already was.

An objective certainty, a confirmation by experience, can never be
obtained. Not only is the future, as has often been remarked, ever
equally uncertain, but when one contemplates nature in order to find
confirmation of God's omnipotence and wisdom, one finds alongside
traces in that direction also traces in the opposite direction. From
objective uncertainty, then, we never escape---and since it is
nevertheless a matter of life and death, this objective uncertainty is
embraced with infinite passion. Truth can be held fast only in the form
of subjectivity, as the object of personal feeling and passion.

\medskip

\noindent\textbf{6.} After Kierkegaard has thus developed the claim that
truth can be grasped and held fast only in a subjective, personal
manner---that truth is subjectivity---he reverses the proposition and
maintains: subjectivity is the truth. For only that which is grasped
with subjective energy and passion can be the truth. If there are two
people, one of whom prays in personal untruth to the true God, while
the other ``with the whole passion of infinity'' prays to an idol, then
the former in reality prays to an idol, while the latter truly prays to
God. What matters first and foremost is a \emph{how}, not a \emph{what};
personal inwardness and energy, not the content to which this inwardness
attaches itself. And with objective uncertainty, with the contradiction
that the content presents to the one who must hold fast to it in
personal existence, inwardness and passion increase. Truth is a venture.
To the proposition that subjectivity is the truth corresponds the other
proposition that truth (objectively considered) is the Paradox.

With this famous proposition of his, that subjectivity is the truth,
Kierkegaard has sharpened Lessing's proposition that only in eternal
striving can one relate oneself to truth. The entire value is shifted
from the objective (dogmatic) to the subjective (psychological) pole.
The meaning is really this: only when truth is appropriated and held
fast in personal feeling and striving does it have value. It is the
personal nutritive value of truth that is decisive. The standard lies in
the movement that is awakened in the inner life of the personality.

In fact, Søren Kierkegaard has in his own way stated the same thing that
Ludwig Feuerbach had maintained a few years earlier in \emph{The Essence
of Christianity}: that theology is psychology. (\emph{The Essence
of Christianity} appeared in 1841, the \emph{Concluding Unscientific
Postscript} in 1846. In this book, and already in \emph{Stages on
Life's Way}, Kierkegaard expressly refers to Feuerbach's interpretation
of Christianity.) Only through an inconsistency can Kierkegaard demand
a definite content, a definite objective faith, of subjectivity: the
main thing is that it be in a state of truth-striving; what it relates
itself to must be a matter of indifference. It must be the fullness and
height of the personal life that matters, and it must provisionally
remain an open question whether a certain definite dogmatic content of
faith is necessary for the development of such a life.

But he does not draw the full consequence of his proposition. That would
have led him to proclaim the principle of the free personality. He does
not abandon the assumption of a standard for truth different from the
personality itself, which for him is given in ecclesiastical dogmatics.
And yet---on closer inspection (which we shall undertake when we have
come to know his ethics) it will become apparent that this dogmatic
guideline is not the ultimate one. He is too much a thinker to be able
to rest content with it. His ultimate standard is of a formal nature:
the relation to the dogmatic-ecclesiastical content of faith is for him
the highest form of personal life, because in this relation the
strongest contradiction, the greatest suffering and passion, comes forth
in subjectivity. The last and decisive objection against him will then
prove to be that he has torn the personality loose from the real, natural
and social context in which alone its life can acquire a positive
content. For only in the free subjectivity, struggling with reality and
learning from reality, can truth dwell---not in the subjectivity that
through formal, ascetic exercises torments itself into obedience to a
content that is contrary to nature and hostile to life.---But only
through an examination of Kierkegaard's ethics can this point be more
clearly illuminated.

\bigskip

What is great about Søren Kierkegaard's epistemology is the seriousness
with which he grasps the connection of thought with personal life. There
is something Greek in him. He wants to return from abstract thinking and
contemplation to personal, existential thinking.

Our ideas all too easily become soldiers in peacetime. Søren Kierkegaard
has led ideas forth from the comfort of contemplation to the strife and
unrest of the struggle of life, and has pointed out the arena where the
validity of ``essential'' knowledge must stand its test. Had the
dogmatic impediment not been there, his thoughts would have attained a
more directly universal significance; as it stands, a translation, a
process of sifting, must take place before they can come to play the
role for humanity that is rightfully theirs.

% ------------------------------------------------------------
\section*{B.\quad Ethics}
\label{sec:4B}
\markright{IV. Søren Kierkegaard's Philosophy --- B. Ethics}

% [text to be added: section B intro, pp. 70]

\subsection*{a.\quad The Leap}
\label{sec:4Ba}
% [text to be added: pp. 70--82]

\subsection*{b.\quad The Stages}
\label{sec:4Bb}
% [text to be added: section b intro, pp. 82]

\subsubsection*{$\alpha$.\quad The Aesthetic View of Life}
\label{sec:4Bb-alpha}
% [text to be added: pp. 83--91]

\subsubsection*{$\beta$.\quad The Ethical View of Life}
\label{sec:4Bb-beta}
% [text to be added: pp. 92--109]

\subsubsection*{$\gamma$.\quad The Religious View of Life}
\label{sec:4Bb-gamma}
% [text to be added: pp. 109--119]

\subsection*{c.\quad The Standard}
\label{sec:4Bc}
% [text to be added: pp. 119--126]

% ============================================================
%  V.  SØREN KIERKEGAARD AND CHRISTIANITY  (pp. 127--149)
% ============================================================
\bookchapter{V.\quad Søren Kierkegaard and Christianity}
\label{ch:5}

\section*{A.\quad Personal Breakthrough}
\label{sec:5A}
\markright{V. Søren Kierkegaard and Christianity --- A. Personal Breakthrough}
% [text to be added: pp. 127--143]

\section*{B.\quad The Last Word}
\label{sec:5B}
\markright{V. Søren Kierkegaard and Christianity --- B. The Last Word}
% [text to be added: pp. 143--149]

% ============================================================
%  CONCLUSION  (pp. 150--159)
% ============================================================
\bookchapter{Conclusion}
\label{ch:slutning}

% [text to be added: pp. 150--159]

\end{document}
